% !TeX root = ../../Book5.tex
% 纲要标题：## 第7章　滤过、Artin–Rees 与完备化
% 纲要位置：几何学计划纲要.md，第 5258 行。
% 正文状态：初稿；尚未正式审核。

% 纲要内容开始
% 
% > 本章目标与证明范围：以形式邻域为何保持有限模正合列为主问题，连续完成 Artin–Rees、张量表示和完备化忠实平坦性。逆极限的一般非正合性必须在证明中实际处理。
% 
% 纲要内容结束

\chapter{滤过、Artin–Rees 与完备化}
\label{g5:ch:07}
\BookChapterNumber{g5:ch:07}

\begin{chapterabstract}
	本章由理想的幂构造局部拓扑与形式邻域。Rees 环把无限多个滤过层放入一个有限型分次对象，Artin–Rees 引理据此比较子模的诱导滤过和本来滤过。通过可数逆系统的提升引理，有限模的完备化保持短正合列，并等于与完备化环张量。随后证明完备化环的 Noether 性、平坦性及局部情形的忠实性，比较维数、伴随分次和有限阶长度。最后用代数幂级数与形式方程说明相容截断、单根提升、收敛和代数化之间的联系。
\end{chapterabstract}

\begin{learninggoals}
	\begin{itemize}
		\item 从理想幂及其商写出滤过、拓扑、Rees 环和伴随分次，并计算有限阶长度的增长。
		\item 完整证明 Artin–Rees 引理，说明它如何控制子模滤过，并推出 Krull 交定理。
		\item 通过稳定像或满射核系统选择相容提升，指出逆极限可能失去右正合性的具体位置。
		\item 建立有限模完备化的正合性、张量表示和完备化环的 Noether 性，再推出忠实平坦及局部检测。
		\item 比较完备化前后的维数和切锥，计算形式方程的逐阶解，并区分数值收敛、代数关系和指定精度的逼近。
	\end{itemize}
\end{learninggoals}

设 \(A=k[t]\)，考虑短正合列
\[
	0\longrightarrow(t)\longrightarrow A\longrightarrow k\longrightarrow0.
\]
若对每一项直接模去 \(t^n\)，第一张映射变为
\[
	(t)/(t^{n+1})\longrightarrow A/(t^n).
\]
其中 \(t^n\) 的类非零，却映到零。于是，“在每个有限阶保持原来的短正合列”已经失败。但把所有阶合成形式级数后，我们期待的列
\[
	0\longrightarrow t\,k[[t]]\longrightarrow k[[t]]
	\longrightarrow k\longrightarrow0
\]
确实正合。有限阶的核如何在极限中得到正确处理，是本章的起点。

\ProseParagraphBreak

对子模 \(N\subseteq M\)，真正嵌入 \(M/I^nM\) 的对象是
\(N/(N\cap I^nM)\)，而按 \(N\) 自己的理想进滤过取商，得到的是 \(N/I^nN\)。Artin–Rees 引理保证在 Noether 环的有限模上，这两组子模只相差一个固定的阶数偏移，从而给出同一个完备化。另一个问题来自逆极限：每一阶有原像，并不自动提供相容的原像。本章将这两项问题分别解决，再把所得正合性与第六章的平坦判据连接起来。

\ProseParagraphBreak

本章使用第二章的 Noether 模、Nakayama 引理和张量积，第三章的长度、参数系及高度定理，第五章的伴随分次和形式方程，以及第六章的平坦与忠实平坦判据。主要参考为\cite[第8节，第55—63页；第13节，第93—98页]{Matsumura1989}。Hilbert–Samuel 部分建立当前所需的最终多项式性和增长上界；局部维数的完备化比较则由参数系与素链提升直接证明。

% !TeX root = ../../Book5.tex
% 纲要标题：### 7.1 滤过及伴随分次
% 纲要位置：几何学计划纲要.md，第 5262 行。
% 正文状态：初稿；尚未正式审核。

% 纲要内容开始
% 
% * 理想进拓扑和分离性
% * Rees 环及有限生成
% * 伴随分次与切锥
% * Hilbert–Samuel 函数的所需增长工具
% 
% 纲要内容结束

\section{滤过及伴随分次}
\label{g5:sec:0701}

第五章在局部环中逐层考察 \(\mathfrak m^n/\mathfrak m^{n+1}\)，用最低次项描述切锥。这些层也可以组成一种逼近方式：两个元素模 \(\mathfrak m^n\) 相同，就表示在前 \(n-1\) 阶无法区别它们。本节把这个想法推广到任意理想 \(I\)，先说明“越来越高阶”给出的拓扑，再把全部理想幂放进一个分次环。这样既能利用 Noether 性同时控制无限多层，也能把有限阶邻域的长度组织成最终稳定的多项式规律。

\subsection{理想进拓扑和分离性}
\label{g5:subsec:0701:01}

设 \(A\) 为交换含幺环，\(I\subseteq A\) 为理想，\(M\) 为 \(A\) 模。子模链
\[
	M=I^0M\supseteq IM\supseteq I^2M\supseteq\cdots
\]
把模元素按能被多少个 \(I\) 中的系数乘出来进行分层。若 \(m-m'\in I^nM\)，则它们在截断 \(M/I^nM\) 中相同；随着 \(n\) 增大，允许忽略的误差越来越小。

\begin{definition}\label{g5:def:adic-topology}
	设 \(A\) 为交换含幺环，\(I\subseteq A\) 为理想，\(M\) 为 \(A\) 模。将
	\[
		M=F^0M\supseteq F^1M\supseteq F^2M\supseteq\cdots
	\]
	这样的递降子模链称为 \(M\) 的一个递降\term[lvguo]{滤过}[filtration]。若还满足 \(IF^nM\subseteq F^{n+1}M\)，称它为 \(I\) 滤过。特别地，\(F^nM=I^nM\) 称为 \(I\) 进滤过。

	以全部陪集 \(m+I^nM\)、\(m\in M,\ n\geq0\)，为开集基所定义的拓扑，称为 \(M\) 的 \(I\) 进拓扑，也称\term[lixiangjintuopu]{理想进拓扑}[adic topology]。若任意两个不同元素都有不相交的开邻域，就称这个拓扑为分离的。
\end{definition}

这些陪集确实构成拓扑基：它们覆盖 \(M\)；若某点同时落在两个陪集中，取两个指数中较大者，对应的更小陪集便包含在二者交中。每个 \(I^nM\) 既开又闭，因为它的补集是其余陪集的并。加法和取负连续；若同时给 \(A\) 配置 \(I\) 进拓扑，则标量乘法也连续。事实上，当 \(a-a_0\in I^n\)、\(m-m_0\in I^nM\) 时，
\[
	am-a_0m_0=a(m-m_0)+(a-a_0)m_0\in I^nM.
\]
这里表达的只是误差落在给定理想幂中，没有使用数值意义的距离。

\begin{proposition}\label{g5:prop:adic-closure}
	设 \(A\) 为交换含幺环，\(I\subseteq A\) 为理想，\(M\) 为 \(A\) 模，\(N\subseteq M\) 为子模。则：
	\begin{enumerate}
		\item \(M\) 的 \(I\) 进拓扑分离，当且仅当 \(\bigcap_{n\geq0}I^nM=0\)；
		\item \(N\) 在 \(M\) 中的闭包为
		\[
			\overline N=\bigcap_{n\geq0}(N+I^nM);
		\]
		\item 商拓扑是 \(M/N\) 自身的 \(I\) 进拓扑；它分离，当且仅当 \(N\) 在 \(M\) 中闭。
	\end{enumerate}
\end{proposition}
\begin{proof}
	若 \(m-m'\notin I^nM\)，两个陪集 \(m+I^nM\) 与 \(m'+I^nM\) 不相交，因为相交会使其差属于 \(I^nM\)。因此交为零时，任意两个不同点都能由某个指数区别。反过来，若 \(0\ne x\in\bigcap_nI^nM\)，则 \(0\) 与 \(x\) 在每一层都属于相同陪集，任意开邻域都无法把它们分开。这证明第一项。

	一个 \(m\in M\) 属于 \(\overline N\)，恰好表示每个基本邻域 \(m+I^nM\) 都与 \(N\) 相交。这又等价于 \(m\in N+I^nM\) 对每个 \(n\) 成立，得到第二项。

	商映射把 \(I^nM\) 送到
	\[
		(I^nM+N)/N=I^n(M/N).
	\]
	它把开陪集送到这些子模的陪集，而且这些陪集的逆像也开，故商拓扑正是所述理想进拓扑。根据第一项，商分离等价于
	\(\bigcap_n(N+I^nM)/N=0\)，再由第二项等价于 \(\overline N=N\)。
\end{proof}

例如 \(A=k[x]\)、\(I=(x)\) 时，非零多项式不可能被任意高次的 \(x\) 整除，所以 \(A\) 分离。但若改取 \(M=k[x,x^{-1}]\)，则 \(x\) 在 \(M\) 上可逆，所有 \(I^nM\) 都等于 \(M\)，其拓扑不分离。另一个有限模例子是 \(A=k\times k\)、\(I=k\times0\)、\(M=A\)：因 \(I^2=I\ne0\)，所有高次截断都忽略了第一因子。分离性因此需要实际检查幂的交，不能仅由环 Noether 或模有限生成单独推出。

\ProseParagraphBreak

任意 \(A\) 线性映射 \(f:M\to M'\) 满足 \(f(I^nM)\subseteq I^nM'\)，所以关于理想进拓扑连续。对子模 \(N\subseteq M\)，则要区分两种滤过：自身的 \(I^nN\) 与从 \(M\) 诱导的 \(N\cap I^nM\)。总有
\[
	I^nN\subseteq N\cap I^nM,
\]
所以自身的理想进拓扑至少与诱导拓扑同样精细。下一节将分析何时两种滤过可以相互控制。上述拓扑、闭包及商拓扑的讨论可对照\cite[第8节，第55--56页]{Matsumura1989}。

\begin{proposition}\label{g5:prop:adic-equivalent-powers}
	设 \(A\) 为交换含幺环，\(I,J\subseteq A\) 为有限生成理想，且 \(\sqrt I=\sqrt J\)。则存在正整数 \(a,b\)，使
	\[
		I^a\subseteq J,\qquad J^b\subseteq I.
	\]
	因此对每个 \(A\) 模 \(M\)，其 \(I\) 进拓扑与 \(J\) 进拓扑相同。
\end{proposition}
\begin{proof}
	写 \(I=(x_1,\ldots,x_r)\)。对每个 \(i\)，由 \(x_i\in\sqrt J\) 选 \(e_i\geq1\)，使 \(x_i^{e_i}\in J\)。若 \(r>0\)，取
	\(a=1+\sum_i(e_i-1)\)。任意总次数为 \(a\) 的生成元乘积中，必有某个 \(x_i\) 出现至少 \(e_i\) 次，所以 \(I^a\subseteq J\)；\(I=0\) 时直接取 \(a=1\)。交换 \(I,J\) 得到另一个包含。

	于是对每个 \(n\)，有 \(I^{an}M\subseteq J^nM\) 及 \(J^{bn}M\subseteq I^nM\)。每一族基本邻域都包含另一族中的某个基本邻域，故所定义的拓扑相同。
\end{proof}

特别地，在 Noether 局部环 \((A,\mathfrak m)\) 中，根式为 \(\mathfrak m\) 的真理想 \(I\) 与 \(\mathfrak m\) 给出同一拓扑。不过，它们在固定指数处的截断长度可以不同，因为相互包含通常需要改变指数。

\subsection{Rees 环及有限生成}
\label{g5:subsec:0701:02}

理想幂形成无限链，但每个 \(I^n\) 都由 \(I\) 中元素的 \(n\) 重乘积生成。要同时使用这些乘法关系，可以给第 \(n\) 层加上一个记录次数的符号 \(t^n\)，再将各层放在同一个环内。这样，关于所有高次幂的问题便能利用一个环的有限生成性处理。

\begin{definition}\label{g5:def:rees-algebra}
	设 \(A\) 为交换含幺环，\(I\subseteq A\) 为理想，\(t\) 为不定元。子环
	\[
		\mathcal R_I(A)\coloneqq\bigoplus_{n\geq0}I^nt^n
		\subseteq A[t]
	\]
	称为 \(I\) 的\term[Reeshuan]{Rees 环}[Rees algebra]。\footnote{David Rees（1918--2013），英国数学家，研究半群与交换代数，在理想约化和理想的渐近理论中作出重要贡献。}其中 \(I^nt^n\) 为次数 \(n\) 的齐次部分。
\end{definition}
\symidx{\(\mathcal R_I(A)\)：理想 \(I\) 的 Rees 环}

这里用的是直和，每个元素只有有限多个非零齐次部分。它确实是子环，因为
\[
	(I^it^i)(I^jt^j)\subseteq I^{i+j}t^{i+j},
\]
且次数零部分为含单位的 \(A\)。次数只由 \(t\) 的幂记录；即使 \(A\) 本身是多项式环，\(A\) 中的元素在这份分次中也都位于次数零。

\begin{proposition}\label{g5:prop:rees-noetherian}
	设 \(A\) 为交换含幺环，\(I=(a_1,\ldots,a_r)\) 为有限生成理想。则
	\[
		\mathcal R_I(A)=A[a_1t,\ldots,a_rt],
	\]
	从而它是有限型 \(A\) 代数。若 \(A\) 为 Noether 环，则 \(\mathcal R_I(A)\) 为 Noether 环。
\end{proposition}
\begin{proof}
	右侧的每个生成元属于 \(\mathcal R_I(A)\)，所以右侧包含于左侧。反过来，\(I^n\) 由所有
	\(a_1^{\alpha_1}\cdots a_r^{\alpha_r}\)、\(\sum_i\alpha_i=n\)，生成。因此每个 \(ct^n\)、\(c\in I^n\)，都是相应乘积
	\((a_1t)^{\alpha_1}\cdots(a_rt)^{\alpha_r}\)
	的 \(A\) 线性组合，得到另一包含。于是同态
	\[
		A[T_1,\ldots,T_r]\longrightarrow\mathcal R_I(A),
		\qquad T_i\longmapsto a_it
	\]
	满射。若 \(A\) Noether，Hilbert 基定理使左侧 Noether，其商环也 Noether。
\end{proof}

例如 \(A=k[x,y]\)、\(I=(x,y)\) 时，Rees 环由次数一元素 \(xt,yt\) 在 \(A\) 上生成。这两个元素满足
\[
	y(xt)-x(yt)=0,
\]
所以有限型并不表示它们在 \(A\) 上代数独立。若 \(I=(a)\) 且 \(a\) 不是零因子，则 \(\mathcal R_I(A)\cong A[T]\)，同构将 \(T\) 送到 \(at\)：一个关系
\(\sum_nc_n(at)^n=0\) 逐次给出 \(c_na^n=0\)，因此全部 \(c_n=0\)。零因子情形则会保留额外关系，例如 \(A=k[\varepsilon]/(\varepsilon^2)\)、\(I=(\varepsilon)\) 中有 \(\varepsilon(\varepsilon t)=0\) 及 \((\varepsilon t)^2=0\)。

\ProseParagraphBreak

上述非负分次环在\cite[第15节，第4部分，第120页]{Matsumura1989}中记为 \(R_+\)。该书还把包含负次数的扩充构造称为 Rees 环；本章的 \(\mathcal R_I(A)\) 始终采用定义中的非负次数约定。它的有限生成性将使下一节能够同时控制子模与各次 \(I^nM\) 的交。

\subsection{伴随分次与切锥}
\label{g5:subsec:0701:03}

Rees 环保留了每个 \(I^n\) 的全部元素。若只保留第 \(n\) 层相对于下一层的差别，就得到伴随分次。第五章已经在局部环的极大理想处定义了这一构造；现在把同一乘法推广到一般理想及模。

\begin{definition}\label{g5:def:ideal-associated-graded}
	设 \(A\) 为交换含幺环，\(I\subseteq A\) 为理想，\(M\) 为 \(A\) 模。定义
	\[
		\operatorname{gr}_I(A)=\bigoplus_{n\geq0}I^n/I^{n+1},
		\qquad
		\operatorname{gr}_I(M)=\bigoplus_{n\geq0}I^nM/I^{n+1}M.
	\]
	第一项称为 \(A\) 关于 \(I\) 的伴随分次环，第二项称为 \(M\) 关于 \(I\) 的\term[bansuifencimo]{伴随分次模}[associated graded module]。其乘法与作用由
	\[
		[a]_i[b]_j=[ab]_{i+j},\qquad
		[a]_i[m]_j=[am]_{i+j}
	\]
	给出，其中方括号下标表示所处次数。
\end{definition}
\symidx{\(\operatorname{gr}_I(A),\operatorname{gr}_I(M)\)：理想进伴随分次环与模}

若改变 \(a\in I^i\) 的代表，变化属于 \(I^{i+1}\)，乘 \(I^j\) 或 \(I^jM\) 后便落入应除去的下一层；改变另一代表时同理。于是乘法与模作用良定义。零次部分为 \(A/I\) 及 \(M/IM\)，每个齐次部分都是 \(A/I\) 模。这一说明也表明，计算各层不需要在 \(A\) 内选出一个与 \(A/I\) 同构的子环。

\begin{proposition}\label{g5:prop:rees-associated-graded}
	设 \(A\) 为交换含幺环，\(I\subseteq A\) 为理想。则有自然的分次环同构
	\[
		\mathcal R_I(A)/I\mathcal R_I(A)
		\cong\operatorname{gr}_I(A),\qquad
		at^n\longmapsto a+I^{n+1}.
	\]
	若 \(I=(a_1,\ldots,a_r)\) 有限生成，则 \(\operatorname{gr}_I(A)\) 是由 \(a_i+I^2\) 这些次数一元素生成的有限型 \(A/I\) 代数。若再设 \(A\) Noether，则 \(\operatorname{gr}_I(A)\) Noether。若 \(M\) 是有限生成 \(A\) 模，则 \(\operatorname{gr}_I(M)\) 是有限生成 \(\operatorname{gr}_I(A)\) 模，由 \(M/IM\) 中有限个次数零元素生成。
\end{proposition}
\begin{proof}
	逐次取商的映射保持乘法并满射。其次数 \(n\) 的核是 \(I^{n+1}t^n\)，而
	\[
		I\mathcal R_I(A)
		=\bigoplus_{n\geq0}I^{n+1}t^n.
	\]
	所以核正是这个理想，得到第一断言。次数一生成的结论既可由 Rees 环的生成元取商得到，也可直接由 \(I^n\) 的生成乘积看出；具体地，有分次满射
	\[
		(A/I)[X_1,\ldots,X_r]\twoheadrightarrow\operatorname{gr}_I(A),
		\qquad X_i\longmapsto a_i+I^2.
	\]
	若 \(A\) Noether，左侧 Noether，故其商也 Noether。

	最后写 \(M=Am_1+\cdots+Am_s\)。每个 \(I^nM\) 元素可写成
	\(\sum_j c_jm_j\)、\(c_j\in I^n\)，其在第 \(n\) 层中的类就是
	\(\sum_j[c_j]_n[m_j]_0\)。因此有限个 \([m_j]_0\) 生成整个伴随分次模。
\end{proof}

本命题中的环与模构造可对照\cite[第13节，第93--94、97页]{Matsumura1989}。它没有假设理想进拓扑分离；若一个元素属于所有 \(I^nM\)，它在全部截断中都不可见，这与伴随分次只记录相邻层的差别相符。

\ProseParagraphBreak

取非零 Noether 局部环 \((A,\mathfrak m,\kappa)\)，便恢复第五章的切锥
\[
	\mathrm C(A)=\operatorname{Spec}\bigl(\operatorname{gr}_{\mathfrak m}A\bigr).
\]
由\cref{g5:prop:tangent-cone-in-tangent-space}，次数一元素诱导满射
\[
	\operatorname{Sym}_\kappa(\mathfrak m/\mathfrak m^2)
	\twoheadrightarrow\operatorname{gr}_{\mathfrak m}A.
\]
因此切锥位于切空间内，且其坐标环保留全部次数的乘法关系。当前的推广还容许把 \(I\) 取成一个子空间的方程理想，此时零次环为 \(A/I\)，相邻理想幂描述围绕这个子空间的各层信息。

\ProseParagraphBreak

伴随分次也把截断长度拆成各层长度。在各项长度有限时，由\cref{g5:prop:length-additivity}逐次应用
\[
	0\longrightarrow I^nM/I^{n+1}M
	\longrightarrow M/I^{n+1}M
	\longrightarrow M/I^nM\longrightarrow0
\]
得到
\begin{equation}\label{g5:eq:truncation-graded-length}
	\ell_A(M/I^{n+1}M)
	=\sum_{j=0}^n\ell_A(I^jM/I^{j+1}M).
\end{equation}
右边的每一项可以由分次环的方法计算，累加后便得到有限阶邻域的厚度。

\subsection{Hilbert–Samuel 函数的所需增长工具}
\label{g5:subsec:0701:04}

有限生成的分次模虽然有无限多个齐次部分，新的部分却由有限个生成元及固定的乘法规则产生。为了同时记录它们的长度，先把长度作为形式幂级数的系数。这里不讨论级数在数值处收敛，等式的含义是每个系数分别相等。

\begin{definition}\label{g5:def:hilbert-series}
	设 \(R\) 为交换 Artin 环，\(S=R[X_1,\ldots,X_r]\)，其中每个 \(X_i\) 的次数为 \(1\)。设
	\(E=\bigoplus_{n\geq0}E_n\) 为有限生成分次 \(S\) 模，即 \(S_iE_j\subseteq E_{i+j}\)。它的\term[Hilberthanshu]{Hilbert 函数}[Hilbert function]与\term[Hilbertjishu]{Hilbert 级数}[Hilbert series]分别为
	\[
		h_E(n)=\ell_R(E_n),\qquad
		P_E(z)=\sum_{n\geq0}h_E(n)z^n\in\mathbb Z[[z]].
	\]
\end{definition}

上述长度是有限的。确实，将有限组模生成元分解为各自有限多个齐次部分，所得齐次元素仍生成 \(E\)。若它们的次数为 \(d_1,\ldots,d_s\)，则 \(E_n\) 由有限个
\[
	X^\alpha e_j,\qquad |\alpha|=n-d_j,
\]
在 \(R\) 上生成。Artin 环作为自身的模有有限长度，所以这些有限生成 \(R\) 模也有有限长度。

\ProseParagraphBreak

下面是次数一分次情形的\term[HilbertSerredingli]{Hilbert–Serre 定理}[Hilbert--Serre theorem]。它用变量的乘法映射，把 \(r\) 个变量的问题归约为 \(r-1\) 个变量的问题。\footnote{Jean-Pierre Serre（1926--），法国数学家，在代数拓扑、代数几何和代数数论中作出重要贡献。}

\begin{theorem}[Hilbert–Serre]\label{g5:thm:hilbert-serre}
	设 \(R\) 为交换 Artin 环，\(S=R[X_1,\ldots,X_r]\)，每个变量的次数为 \(1\)，\(E=\bigoplus_{n\geq0}E_n\) 为有限生成分次 \(S\) 模。则存在 \(Q(z)\in\mathbb Z[z]\)，使
	\[
		P_E(z)=\frac{Q(z)}{(1-z)^r}.
	\]
	当 \(r\geq1\) 时，\(h_E(n)\) 对充分大的 \(n\) 等于一个次数至多为 \(r-1\) 的有理系数多项式；允许该多项式为零。当 \(r=0\) 时，\(h_E(n)\) 最终为零。对任意 \(r\geq0\)，累积函数
	\[
		H_E(n)=\sum_{j=0}^n h_E(j)
	\]
	最终等于一个次数至多为 \(r\) 的有理系数多项式，同样允许零多项式。
\end{theorem}
\begin{proof}
	先对 \(r\) 归纳证明级数公式。若 \(r=0\)，则 \(S=R\) 全部位于次数零；有限个齐次生成元只能生成有限多个非零齐次部分，所以 \(P_E(z)\) 本身是多项式。

	设 \(r>0\)，令
	\[
		K=(0:_E X_r),\qquad C=E/X_rE.
	\]
	\(S\) Noether，因而 \(K\) 与 \(C\) 都为有限生成分次 \(S\) 模。又 \(X_rK=X_rC=0\)，所以它们实际上是
	\(S'=R[X_1,\ldots,X_{r-1}]\) 上的有限生成分次模。对每个 \(n\geq0\)，乘以 \(X_r\) 给出正合列
	\[
		0\longrightarrow K_{n-1}\longrightarrow E_{n-1}
		\xrightarrow{\,X_r\,}E_n\longrightarrow C_n\longrightarrow0,
	\]
	其中负次数部分约定为零。长度加性得到
	\[
		h_E(n)-h_E(n-1)=h_C(n)-h_K(n-1).
	\]
	逐次乘以 \(z^n\) 并求和，得
	\[
		(1-z)P_E(z)=P_C(z)-zP_K(z).
	\]
	归纳假设使右边具有分母 \((1-z)^{r-1}\)，于是 \(P_E\) 具有所需形式。这里保留核 \(K\) 十分必要，乘以变量未必单射。

	当 \(r\geq1\) 时，形式恒等式
	\[
		(1-z)^{-r}
		=\sum_{n\geq0}\binom{n+r-1}{r-1}z^n
	\]
	来自 \(r\) 份几何级数的乘积：系数计算的是 \(n\) 分成 \(r\) 个非负整数之和的方法数。写 \(Q(z)=\sum_{j=0}^q c_jz^j\)，则对 \(n\geq q\)，有
	\[
		h_E(n)=\sum_{j=0}^q c_j
		\binom{n-j+r-1}{r-1}.
	\]
	每个二项式系数在这里都是关于 \(n\) 的次数 \(r-1\) 多项式，所以右侧为次数至多 \(r-1\) 的多项式。若 \(Q=0\)，结论直接成立。

	最后，逐项累加等价于再乘以 \((1-z)^{-1}\)，即
	\[
		\sum_{n\geq0}H_E(n)z^n
		=\frac{P_E(z)}{1-z}
		=\frac{Q(z)}{(1-z)^{r+1}}.
	\]
	相同的系数计算给出
	\[
		H_E(n)=\sum_{j=0}^q c_j\binom{n-j+r}{r}
		\qquad(n\geq q),
	\]
	其次数至多为 \(r\)。这个公式也包含 \(r=0\) 的最终常数情形。
\end{proof}

最终与 \(h_E(n)\) 相等的多项式称为分次模的\term[Hilbertduoxiangshi]{Hilbert 多项式}[Hilbert polynomial]；它唯一，因为两个多项式若在所有充分大的整数上相同，其差有无限多个根，必为零。定理及系数展开可对照\cite[第13节，定理13.2及其推论，第94--96页]{Matsumura1989}。

\ProseParagraphBreak

变量次数为 \(1\) 是最终单个多项式描述的关键条件。若给 \(k[u]\) 配置 \(\deg u=2\)，则其齐次部分维数在偶数次为 \(1\)、奇数次为 \(0\)，级数为 \((1-z^2)^{-1}\)。这个函数不能最终等于一个多项式：这样的多项式会在无限多个奇数处为零，因而恒为零，却又必须在充分大的偶数处等于 \(1\)。理想进伴随分次由次数一生成，恰好适用上述定理。

\begin{definition}\label{g5:def:hilbert-samuel-function}
	设 \((A,\mathfrak m)\) 为非零 Noether 局部环，\(I\subseteq A\) 为满足 \(\sqrt I=\mathfrak m\) 的真理想，\(M\) 为有限生成 \(A\) 模。定义其关于 \(I\) 的\term[HilbertSamuelhanshu]{Hilbert–Samuel 函数}[Hilbert--Samuel function]为
	\[
		H_{I,M}(n)\coloneqq\ell_A(M/I^{n+1}M),\qquad n\geq0.
	\]
	当 \(M=A\) 时简写为 \(H_{I,A}\)；当 \(I=\mathfrak m\) 时，它记录闭点各阶截断的长度。\footnote{Pierre Samuel（1921--2009），法国数学家，研究交换代数及其在代数几何中的应用，特别是局部环与重数理论。}
\end{definition}
\symidx{\(H_{I,M}(n)\)：长度 \(\ell_A(M/I^{n+1}M)\)}

由第三章的零维与有限长度结论，\(A/I^{n+1}\) 是 Artin 局部环，因为其根式仍为 \(\mathfrak m\)。所以有限生成模 \(M/I^{n+1}M\) 确有有限长度。这里用 \(n+1\) 次幂，是为了让 \(n=0\) 正好对应零次分次部分 \(M/IM\)。

\begin{proposition}\label{g5:prop:hilbert-samuel-polynomial}
	设 \((A,\mathfrak m)\) 为非零 Noether 局部环，\(I=(a_1,\ldots,a_r)\) 为根式等于 \(\mathfrak m\) 的真理想，\(M\) 为有限生成 \(A\) 模。则 \(H_{I,M}(n)\) 最终等于一个有理系数多项式 \(P_{I,M}(n)\)，其次数至多为 \(r\)。若 \(M\ne0\)，则这个多项式非零且首项系数为正。若 \(M\) 可由 \(s\) 个元素生成，则对每个 \(n\geq0\)，都有
	\begin{equation}\label{g5:eq:ideal-generators-growth-bound}
		H_{I,M}(n)\leq
		\ell_A(M/IM)\binom{n+r}{r}
		\leq s\,\ell_A(A/I)\binom{n+r}{r}.
	\end{equation}
\end{proposition}
\begin{proof}
	取 \(R=A/I\)，它是 Artin 环。由\cref{g5:prop:rees-associated-graded}，\(\operatorname{gr}_I(M)\) 是次数一多项式环
	\(R[X_1,\ldots,X_r]\) 上的有限生成分次模。每个齐次部分的 \(R\) 子模与 \(A\) 子模相同，故两种长度一致。根据\eqref{g5:eq:truncation-graded-length}，其累积 Hilbert 函数就是 \(H_{I,M}\)。应用\cref{g5:thm:hilbert-serre}，得到最终多项式性与次数界。

	若 \(M\ne0\)，Nakayama 引理与 \(I\subseteq\mathfrak m\) 保证 \(M/IM\ne0\)，故 \(H_{I,M}(n)\geq1\) 对所有 \(n\) 成立。因此最终多项式不为零，且在充分大的正整数上取正值，其首项系数必为正。

	对每个总次数为 \(j\) 的单项式 \(a^\alpha\)，乘法诱导
	\[
		M/IM\longrightarrow I^jM/I^{j+1}M,\qquad
		m+IM\longmapsto a^\alpha m+I^{j+1}M.
	\]
	这与代表无关，因为 \(a^\alpha IM\subseteq I^{j+1}M\)。合并全部这样的映射得到满射。若 \(r\geq1\)，总共有 \(\binom{j+r-1}{r-1}\) 个单项式，故
	\[
		\ell_A(I^jM/I^{j+1}M)
		\leq\ell_A(M/IM)\binom{j+r-1}{r-1}.
	\]
	从 \(j=0\) 加到 \(n\)，用
	\(\sum_{j=0}^n\binom{j+r-1}{r-1}=\binom{n+r}{r}\)
	得到第一条上界。若 \(r=0\)，则 \(I=0\)，函数恒为 \(\ell_A(M)\)，同一上界仍成立。第二条上界来自满射
	\((A/I)^s\to M/IM\) 及长度加性。
\end{proof}

唯一的最终多项式 \(P_{I,M}\) 称为相应的\term[HilbertSamuelduoxiangshi]{Hilbert–Samuel 多项式}[Hilbert--Samuel polynomial]。累积长度函数与伴随分次的联系见\cite[第13节，第97--98页]{Matsumura1989}。上界中出现的是给定理想的生成元数；借助第三章的参数系，还能把局部环维数放进一个更小的增长上界。

\begin{proposition}\label{g5:prop:parameter-growth-bound}
	设 \((A,\mathfrak m)\) 为非零 Noether 局部环，\(d=\dim A\)，\(I\subseteq A\) 为根式等于 \(\mathfrak m\) 的真理想，\(M\) 为由 \(s\) 个元素生成的有限 \(A\) 模。则存在参数理想
	\[
		\mathfrak q=(z_1,\ldots,z_d)\subseteq I
	\]
	以及有限常数 \(C=s\,\ell_A(A/\mathfrak q)\)，使对每个 \(n\geq0\) 都有
	\[
		H_{I,M}(n)\leq H_{\mathfrak q,M}(n)
		\leq C\binom{n+d}{d}.
	\]
	因此当 \(M\ne0\) 时，\(\deg P_{I,M}\leq d\)。
\end{proposition}
\begin{proof}
	由\cref{g5:thm:system-of-parameters}，先选参数系 \(x_1,\ldots,x_d\)。每个 \(x_i\in\mathfrak m=\sqrt I\)，可取 \(e_i\geq1\) 使 \(z_i=x_i^{e_i}\in I\)。取幂不改变所生成理想的根式，所以 \(\mathfrak q=(z_1,\ldots,z_d)\) 仍为参数理想。若 \(d=0\)，取空参数系，即 \(\mathfrak q=0\)；此时 \(A\) 为 Artin 环。

	由 \(\mathfrak q\subseteq I\)，有 \(\mathfrak q^{n+1}M\subseteq I^{n+1}M\)，于是
	\(M/\mathfrak q^{n+1}M\to M/I^{n+1}M\) 满射，长度给出第一条不等式。把\eqref{g5:eq:ideal-generators-growth-bound}用于由 \(d\) 个元素生成的 \(\mathfrak q\)，便得到第二条不等式。右端为次数 \(d\) 的多项式；\(M\ne0\) 时 \(P_{I,M}\) 的首项系数为正，因此它的次数不可能超过 \(d\)。
\end{proof}

\begin{example}\label{g5:exm:hilbert-samuel-basic-models}
	设 \(k\) 为域。对 \(A=k[x_1,\ldots,x_d]_{(x_1,\ldots,x_d)}\) 及其极大理想 \(\mathfrak m\)，第五章的初始理想计算给出
	\(\operatorname{gr}_{\mathfrak m}A\cong k[X_1,\ldots,X_d]\)。截断中次数至多 \(n\) 的单项式构成基，所以
	\[
		H_{\mathfrak m,A}(n)=\binom{n+d}{d}.
	\]
	当 \(d=0\) 时 \(A=k\)，函数恒为 \(1\)。对任意 DVR \(R\) 及其极大理想 \((\pi)\)，\cref{g5:prop:dvr-length}则给出
	\[
		H_{(\pi),R}(n)=\ell_R(R/\pi^{n+1}R)=n+1.
	\]
	这个长度计算也适用于剩余域不等于某个预先指定基域的情形。

	第五章的尖点与结点模型具有相同的齐次层维数 \(1,2,2,\ldots\)，因此二者都满足
	\[
		H_{\mathfrak m,A}(n)=1+2n.
	\]
	这里结点取\cref{g5:exm:node-tangent-cone}中的特征不为二情形。它们的长度函数相同，切锥环却分别具有二重直线与两个不同方向的结构。相比之下，加厚点 \(k[x,y]/(x,y)^2\) 的长度函数在 \(n=0\) 时为 \(1\)，在 \(n\geq1\) 时恒为 \(3\)。这也说明多项式公式可能只在若干初始层之后才稳定。
\end{example}

最终多项式性使无限多个截断长度具有有限的描述，而参数系给出它们增长速度的上界。长度只计算各层的大小，不能恢复层与层之间的乘法；切锥中的幂零元和不同分支，仍需由伴随分次环本身辨认。
% 纲要条目排为编号小节标题，并进入目录与书签；说明段仍保留于纲要注释中。

% !TeX root = ../../Book5.tex
% 纲要标题：### 7.2 Artin–Rees 与交定理
% 纲要位置：几何学计划纲要.md，第 5269 行。
% 正文状态：初稿；尚未正式审核。

% 纲要内容开始
% 
% * 诱导滤过与本来滤过
% * Rees 模的 Noether 论证
% * Artin–Rees 引理
% * Krull 交定理的局部形式
% 
% 纲要内容结束

\section{Artin–Rees 与交定理}
\label{g5:sec:0702}

设 \(N\) 是模 \(M\) 的子模。若把 \(I^nM\) 看作精度为 \(n\) 时可以忽略的误差，那么 \(N\) 中被忽略的部分是 \(N\cap I^nM\)。但从 \(N\) 自身出发，其 \(I\) 进误差应当是 \(I^nN\)。这两个子模一般不相等，因此在模去 \(I^n\) 后，原来的包含 \(N\hookrightarrow M\) 未必仍为单射。Artin–Rees 引理说明，在 Noether 环的有限模上，两种精度之间至多相差一个固定的整数。将这一结论用于所有阶误差的交，又会得到有限模的分离性及子模的闭性。

\subsection{诱导滤过与本来滤过}
\label{g5:subsec:0702:01}

设 \(A\) 为交换含幺环，\(I\subseteq A\) 为理想，\(M\) 为 \(A\) 模，\(N\subseteq M\) 为子模。对 \(n\geq0\)，记
\[
	F_n=N\cap I^nM,\qquad G_n=I^nN.
\]
它们都是从 \(N\) 开始的降滤过，并满足
\[
	IF_n\subseteq F_{n+1},\qquad IG_n=G_{n+1},\qquad
	G_n\subseteq F_n.
\]
其中 \((F_n)\) 是从 \(M\) 的 \(I\) 进滤过得到的\term[youdaolvguo]{诱导滤过}[induced filtration]；\((G_n)\) 是 \(N\) 自身的 \(I\) 进滤过，本节也称它为本来滤过。以 \(F_n\) 为零点邻域得到的拓扑正是 \(N\subseteq M\) 的子空间拓扑，而以 \(G_n\) 为零点邻域得到的是 \(N\) 自身的 \(I\) 进拓扑。由于 \(G_n\subseteq F_n\)，本来拓扑至少和诱导拓扑一样细。

\begin{example}\label{g5:exm:shifted-induced-filtration}
	设 \(A=k[x]\)、\(I=(x)\)、\(M=A\)，取 \(N=x^aA\)，其中 \(a\geq1\)。则
	\[
		I^nM\cap N=x^{\max(n,a)}A,
		\qquad I^nN=x^{n+a}A.
	\]
	例如 \(n\geq a\) 时，前者为 \(x^nA\)，后者为 \(x^{n+a}A\)，并不相等。但把第一种滤过的下标再增加 \(a\)，就有
	\[
		I^{n+a}M\cap N=I^nN\qquad(n\geq0).
	\]
	因此它们给出相同的零点邻域系统，只是对同一精度使用了不同下标。
\end{example}

一般地，两个降滤过 \((F_n)\)、\((G_n)\) 给出同一拓扑，当且仅当它们的零点邻域互相包含到任意深处：对每个 \(n\)，分别存在 \(a,b\) 使 \(F_a\subseteq G_n\)、\(G_b\subseteq F_n\)。这一条件也称两个滤过互相共尾。它不要求 \(F_n=G_n\)，甚至不要求通过固定平移得到相等。

\ProseParagraphBreak

诱导滤过是有限阶商中自然出现的对象。对任意 \(n\geq0\)，有正合列
\begin{equation}\label{g5:eq:induced-filtration-quotients}
	0\longrightarrow N/(N\cap I^nM)
	\longrightarrow M/I^nM
	\longrightarrow (M/N)/I^n(M/N)\longrightarrow0.
\end{equation}
第一张映射由包含诱导，其核已被除去。最后一项等于 \(M/(N+I^nM)\)，故中间项的核正是 \((N+I^nM)/I^nM\)，也就是第一项的像，最后一张映射显然满射。若把第一项直接换成 \(N/I^nN\)，则其到中间项的映射有核
\[
	\frac{N\cap I^nM}{I^nN}.
\]
这个核正好测量两种滤过在第 \(n\) 阶的差别。要让连续提高精度的过程兼容子模，就需要同时控制所有这些核。

\subsection{Rees 模的 Noether 论证}
\label{g5:subsec:0702:02}

上一节用一个分次环同时保存 \(I\) 的所有幂。对模也作同样的构造，就可以把无限多个交 \(N\cap I^nM\) 放进一个模中，转而使用 Noether 条件提供的有限生成性。

\begin{definition}\label{g5:def:rees-module}
	设 \(A\) 为交换含幺环，\(I\subseteq A\) 为理想，\(M\) 为 \(A\) 模。把
	\[
		M[t]=\bigoplus_{n\geq0}Mt^n
	\]
	看作具有通常多项式乘法作用的 \(A[t]\) 模。其中的分次 \(A\) 子模
	\[
		\mathcal R_I(M)=\bigoplus_{n\geq0}I^nM\,t^n
	\]
	称为 \(M\) 的\term[Reesmo]{Rees 模}[Rees module]。它是上一节的 Rees 环
	\(\mathcal R_I(A)=\bigoplus_{n\geq0}I^nt^n\) 上的模，作用由
	\[
		(at^r)(mt^s)=(am)t^{r+s}
		\qquad(a\in I^r,\ m\in I^sM)
	\]
	给出。
\end{definition}

这里 \(t\) 用来记录次数，每个元素仍只有有限个非零齐次分量。若 \(N\subseteq M\)，则
\[
	\mathcal E(N,M)=\bigoplus_{n\geq0}(I^nM\cap N)t^n
\]
是 \(\mathcal R_I(M)\) 的分次 \(\mathcal R_I(A)\) 子模：乘以 \(I^r\) 中的元素，既不会离开 \(N\)，也会把 \(I^nM\) 送入 \(I^{n+r}M\)。它通常大于 \(\mathcal R_I(N)\)，这与刚才两种滤过的差别对应。

\begin{proposition}\label{g5:prop:rees-module-noetherian}
	设 \(A\) 为交换含幺 Noether 环，\(I\subseteq A\) 为理想，\(M\) 为有限生成 \(A\) 模，\(N\subseteq M\) 为子模。则：
	\begin{enumerate}
		\item \(\mathcal R_I(M)\) 是由有限个零次元素生成的 \(\mathcal R_I(A)\) 模，因而是 Noether 模；
		\item \(\mathcal E(N,M)\) 可由有限个齐次元素
		\[
			u_1t^{d_1},\ldots,u_st^{d_s},
			\qquad u_j\in I^{d_j}M\cap N,
		\]
		作为 \(\mathcal R_I(A)\) 模生成。
	\end{enumerate}
\end{proposition}
\begin{proof}
	取 \(M\) 的有限生成元 \(m_1,\ldots,m_r\)。每个 \(v\in I^nM\) 都能写成
	\[
		v=\sum_{j=1}^r a_jm_j,\qquad a_j\in I^n,
	\]
	因为在 \(I^nM\) 的有限和表达式中，再把各模元素用 \(m_j\) 展开即可。因此
	\[
		vt^n=\sum_{j=1}^r(a_jt^n)m_j,
	\]
	说明零次元素 \(m_j\) 生成整个 Rees 模。由\cref{g5:prop:rees-noetherian}，\(\mathcal R_I(A)\) 为 Noether 环；再用\cref{g5:prop:noetherian-finite-modules}，得到 \(\mathcal R_I(M)\) 为 Noether 模。

	已经验证 \(\mathcal E(N,M)\) 是其子模，故它有限生成。取一组有限生成元，将每个生成元拆成有限个齐次分量。因为 \(\mathcal E(N,M)\) 是直和定义的分次子模，这些分量仍属于 \(\mathcal E(N,M)\)。把所有分量合在一起，仍为有限组；它们生成的子模包含原生成元，故生成整个 \(\mathcal E(N,M)\)。这便得到陈述中的齐次生成组。\(N=0\) 时允许空生成组。
\end{proof}

Noether 条件在这里控制的不是某一个固定的交，而是所有阶的交组成的分次子模。只分别知道每个 \(I^nM\cap N\) 有限生成，还不能保证生成元出现的次数有统一上界；上面的命题正好给出这个上界。

\subsection{Artin–Rees 引理}
\label{g5:subsec:0702:03}

设上一命题中齐次生成元的次数不超过 \(c\)。当 \(n\geq c\) 时，一个 \(n\) 次元素只能由这些较低次数的生成元乘上相应次数的系数得到。把系数中共同的 \(I^{n-c}\) 部分提出，就能把全部高阶交统一表示为第 \(c\) 阶交的倍数。

\begin{theorem}[Artin–Rees 引理]\label{g5:thm:artin-rees}
	设 \(A\) 为交换含幺 Noether 环，\(I\subseteq A\) 为理想，\(M\) 为有限生成 \(A\) 模，\(N\subseteq M\) 为子模。存在整数 \(c\geq0\)，使对每个 \(n\geq c\) 都有
	\begin{equation}\label{g5:eq:artin-rees}
		I^nM\cap N=I^{n-c}(I^cM\cap N).
	\end{equation}
\end{theorem}
\begin{proof}
	若 \(N=0\)，取 \(c=0\) 即可。否则由\cref{g5:prop:rees-module-noetherian}，选定 \(\mathcal E(N,M)\) 的有限齐次生成元 \(u_jt^{d_j}\)，并令 \(c=\max_jd_j\)。对 \(v\in I^nM\cap N\)，存在 \(b_j(t)\in\mathcal R_I(A)\) 使
	\[
		vt^n=\sum_{j=1}^s b_j(t)u_jt^{d_j}.
	\]
	等式两边取 \(n\) 次齐次分量，只有 \(b_j(t)\) 的 \(n-d_j\) 次部分能够留下。若 \(n\geq c\)，便得到
	\[
		I^nM\cap N=\sum_{j=1}^s I^{n-d_j}u_j.
	\]
	这里从右到左的包含也直接成立，因为 \(u_j\in I^{d_j}M\cap N\)。特别地，在 \(n=c\) 时有
	\[
		I^cM\cap N=\sum_{j=1}^s I^{c-d_j}u_j.
	\]
	对 \(n\geq c\)，利用理想幂的乘法，
	\[
		\begin{aligned}
			I^nM\cap N
			&=\sum_{j=1}^s I^{n-c}I^{c-d_j}u_j\\
			&=I^{n-c}\left(\sum_{j=1}^s I^{c-d_j}u_j\right)
			=I^{n-c}(I^cM\cap N).
		\end{aligned}
	\]
	这就是所需等式。
\end{proof}

这一结论称为\term[ArtinReesyinli]{Artin–Rees 引理}[Artin–Rees lemma]。证明把无限多阶的关系归结为有限齐次生成元的次数；同一 Noether 论证也见\cite[第8节，定理8.5，第59页]{Matsumura1989}。常数 \(c\) 可以依赖于 \(I,M,N\)，并不是只由底环决定的统一常数。

\begin{example}\label{g5:exm:artin-rees-shift}
	在\cref{g5:exm:shifted-induced-filtration}中，取 \(c=a\) 就有
	\[
		I^nM\cap N=x^nA
		=I^{n-a}(I^aM\cap N)\qquad(n\geq a).
	\]
	而 \(c<a\) 不可能满足引理的全部等式：在 \(n=a\) 时，左侧为 \(x^aA\)，右侧却为 \(x^{2a-c}A\)，严格更小。因此即使固定 \(A=k[x]\)、\(I=(x)\)、\(M=A\)，当子模 \(x^aA\) 改变时，所需的 \(c\) 仍可任意增大。
\end{example}

\begin{corollary}\label{g5:cor:induced-adic-topology}
	设 \(A\) 为交换含幺 Noether 环，\(I\subseteq A\) 为理想，\(M\) 为有限生成 \(A\) 模，\(N\subseteq M\) 为子模。存在 \(c\geq0\)，使对所有 \(n\geq0\) 都有
	\[
		I^{n+c}M\cap N\subseteq I^nN\subseteq I^nM\cap N.
	\]
	所以 \(N\) 的本来 \(I\) 进滤过与从 \(M\) 诱导的滤过互相共尾，给出同一拓扑。
\end{corollary}
\begin{proof}
	取 Artin–Rees 引理中的 \(c\)。对 \(n+c\geq c\) 使用等式，得到
	\[
		I^{n+c}M\cap N
		=I^n(I^cM\cap N)\subseteq I^nN.
	\]
	另一包含由 \(N\subseteq M\) 直接得到。这两侧包含说明，两种零点邻域系统中的每个邻域，都包含另一系统中足够深的一个邻域，因而拓扑相同。
\end{proof}

这正是\cite[第8节，定理8.6，第59—60页]{Matsumura1989}中的拓扑结论。它保证有限阶正合列\eqref{g5:eq:induced-filtration-quotients}左端所用的诱导精度，最终能与 \(N\) 自身的精度比较。拓扑相同并不自动意味着 \(N\) 是 \(M\) 中的闭集；闭性还要检查商模是否分离。

\subsection{Krull 交定理的局部形式}
\label{g5:subsec:0702:04}

一个元素若属于每个 \(I^nM\)，就在所有有限阶商 \(M/I^nM\) 中都看不见。令
\[
	L=\bigcap_{n\geq0}I^nM.
\]
要证明 \(I\) 进拓扑分离，就要证明 \(L=0\)。Artin–Rees 先给出对任意理想都成立的关系 \(IL=L\)，再由 Nakayama 的行列式形式决定这个交何时必须消失。

\begin{proposition}\label{g5:prop:krull-intersection-annihilator}
	设 \(A\) 为交换含幺 Noether 环，\(I\subseteq A\) 为理想，\(M\) 为有限生成 \(A\) 模，\(L=\bigcap_{n\geq0}I^nM\)。则
	\[
		IL=L,
	\]
	并且存在 \(a\in I\)，使 \((1-a)L=0\)。
\end{proposition}
\begin{proof}
	\(L\) 是有限模 \(M\) 的子模，由 Noether 条件知它有限生成。对包含 \(L\subseteq M\) 应用 Artin–Rees 引理，取相应常数 \(c\)。因为 \(L\) 包含于每个 \(I^nM\)，在 \(n=c+1\) 时得到
	\[
		L=I^{c+1}M\cap L
		=I(I^cM\cap L)=IL.
	\]
	对有限模 \(L\) 应用\cref{g5:prop:nakayama-ideal}，便存在 \(a\in I\) 使 \((1-a)L=0\)。这一证明没有把降链 \(I^nM\) 当作最终稳定的链；稳定关系来自 Artin–Rees，而不是 Noether 环的升链条件本身。
\end{proof}

为保证这里的 \(1-a\) 可逆，记
\[
	\operatorname{Jac}(A)=
	\bigcap_{\mathfrak m\text{ 为 }A\text{ 的极大理想}}\mathfrak m,
\]
称为 \(A\) 的\term[Jacobsongen]{Jacobson 根}[Jacobson radical]。
\footnote{Nathan Jacobson（1910--1999），美国数学家，研究环论，并撰写多部代数学著作。}零环中按空交约定取 \(\operatorname{Jac}(0)=0\)。若 \(a\in\operatorname{Jac}(A)\)，则 \(1-a\) 是单位：否则它属于某个极大理想 \(\mathfrak m\)，而 \(a\in\mathfrak m\)，便会得到 \(1\in\mathfrak m\)，矛盾。在局部环 \((A,\mathfrak m)\) 中，\(\operatorname{Jac}(A)=\mathfrak m\)。

\begin{theorem}[Krull 交定理]\label{g5:thm:krull-intersection}
	设 \(A\) 为交换含幺 Noether 环，\(M\) 为有限生成 \(A\) 模，\(I\subseteq\operatorname{Jac}(A)\) 为理想。则
	\[
		\bigcap_{n\geq0}I^nM=0.
	\]
	特别地，若 \((A,\mathfrak m)\) 为 Noether 局部环，\(I\subseteq\mathfrak m\)，则每个有限生成 \(A\) 模的 \(I\) 进拓扑都分离；取 \(I=\mathfrak m\) 得到极大理想进拓扑的分离性。
\end{theorem}
\begin{proof}
	由上个命题，存在 \(a\in I\) 使 \((1-a)L=0\)。因为 \(a\in\operatorname{Jac}(A)\)，\(1-a\) 是单位，故 \(L=0\)。局部情形由 \(\operatorname{Jac}(A)=\mathfrak m\) 立即得到。
\end{proof}

这就是\term[Krulljiaodingli]{Krull 交定理}[Krull intersection theorem]的有限模形式；一般的零化关系及分离性结论，可对照\cite[第8节，定理8.9—8.10，第60页]{Matsumura1989}。若 \(A\) 是 Noether 整环、\(I\ne A\)，即使 \(I\) 未必包含于 Jacobson 根，也有 \(\bigcap_n I^n=0\)：对 \(M=A\) 使用零化关系，因 \(a\in I\) 而有 \(1-a\ne0\)，整环中非零元素不能零化非零元素。

\begin{corollary}\label{g5:cor:adic-submodule-closed}
	设 \(A\) 为交换含幺 Noether 环，\(I\subseteq\operatorname{Jac}(A)\) 为理想，\(M\) 为有限生成 \(A\) 模。则 \(M\) 的每个子模 \(N\) 都在 \(I\) 进拓扑中闭，并且
	\[
		N=\bigcap_{n\geq0}(N+I^nM).
	\]
\end{corollary}
\begin{proof}
	先说明闭包公式。对任意 \(x\in M\)，它属于 \(N\) 的闭包，当且仅当每个邻域 \(x+I^nM\) 都与 \(N\) 相交。这又等价于 \(x\in N+I^nM\) 对每个 \(n\) 成立。因此，对任何模和理想都有
	\[
		\overline N=\bigcap_{n\geq0}(N+I^nM),
		\qquad
		\overline N/N=\bigcap_{n\geq0}I^n(M/N).
	\]
	商模 \(M/N\) 仍有限生成，故 Krull 交定理使右侧为零，得到 \(\overline N=N\)。这也说明，子模闭性准确地对应于商模 \(M/N\) 的 \(I\) 进分离性。
\end{proof}

\begin{example}\label{g5:exm:idempotent-adic-nonseparated}
	设 \(k\) 为域，\(A=k\times k\)，\(I=k\times0\)。这是 Noether 环中的真理想，但 \(I^2=I\)，所以
	\[
		\bigcap_{n\geq0}I^n=I\ne0.
	\]
	其两个极大理想为 \(k\times0\) 和 \(0\times k\)，故 \(\operatorname{Jac}(A)=0\)，本例不满足 \(I\subseteq\operatorname{Jac}(A)\)。所有有限阶商都只留下第二个 \(k\) 分量，第一分量在每一阶已经消失。因而一般 Noether 环上的真理想，并不足以保证其进拓扑分离。
\end{example}

\begin{example}\label{g5:exm:separated-module-dense-submodule}
	设 \(A=k[x]\)、\(I=(x)\)、\(M=A\)、\(N=(x-1)A\)。多项式不可能被任意高次的 \(x\) 整除，故 \(M\) 的 \(I\) 进拓扑分离。然而对每个 \(n\geq1\)，等式
	\[
		1=(1-x)(1+x+\cdots+x^{n-1})+x^n
	\]
	给出 \(N+I^nM=M\)。因此 \(\overline N=M\ne N\)，\(N\) 是一个稠密的真子模。商模 \(M/N\cong k\) 上的 \(x\) 作用为恒等，故其 \(I\) 进滤过每一项都是整个商模，并不分离。

	本例甚至有 \(I^nM\cap N=I^nN\)：若 \(x^nf(x)\) 在 \(x=1\) 处为零，则 \(f(1)=0\)，故 \(f\in(x-1)A\)。所以诱导拓扑与本来拓扑相同，不能代替子模闭性的检验；即使环境模已经分离，也仍须检查商模。
\end{example}

有限生成性同样有实际作用。令 \(A=k[x]_{(x)}\)、\(I=(x)\)，取无限生成的 \(A\) 模 \(M=k(x)\)。因为乘以 \(x\) 在 \(M\) 上可逆，所有 \(I^nM\) 都等于 \(M\)，交非零。这里底环 Noether 且局部，失效的是有限生成条件；\subsecref{g5:subsec:0203:04}中对 Nakayama 引理的反例已经说明这个模不能有限生成。
% 纲要条目排为编号小节标题，并进入目录与书签；说明段仍保留于纲要注释中。

% !TeX root = ../../Book5.tex
% 纲要标题：### 7.3 有限模的完备化
% 纲要位置：几何学计划纲要.md，第 5276 行。
% 正文状态：初稿；尚未正式审核。

% 纲要内容开始
% 
% * 逆系统和 Mittag–Leffler 的最低引理
% * 完备化的正合性
% * 有限模的张量表示
% * 完备化环的 Noether 性
% 
% 纲要内容结束

\section{有限模的完备化}
\label{g5:sec:0703}

形式级数由全部相容截断确定，这一点在\secref{g5:sec:0504}已经用于计算局部方程。对一般有限模，同样可以把 \(M/I^nM\) 拼成一个完整对象。眼前的问题是：短正合列中的元素在每一阶都有原像，能否把这些原像选得彼此相容？逐阶满射本身不足以保证这一点，必须控制不同提升之间的差。

\ProseParagraphBreak

本节先在可数逆系统中建立所需的提升引理，再用 Artin–Rees 比较子模上的两种滤过。由此得到有限模完备化的正合性，并将完备化写成与 \(\widehat A\) 的张量积。最后通过截断商和伴随分次环，证明 \(\widehat A\) 本身仍为 Noether 环。逆极限的基本构造可对照\cite[附录A，第271—273页]{Matsumura1989}，有限模完备化的主要性质见该书第8节。

\subsection{逆系统和 Mittag–Leffler 的最低引理}
\label{g5:subsec:0703:01}

\begin{definition}\label{g5:def:inverse-system}
	设 \(A\) 为交换含幺环。一个以正整数为指标的 \(A\) 模\term[nixitong]{逆系统}[inverse system]由模 \(E_n\) 及转移同态
	\[
		\rho_n:E_{n+1}\longrightarrow E_n\qquad(n\geq1)
	\]
	组成。对 \(m>n\)，记复合为 \(\rho_{m,n}:E_m\to E_n\)，并令 \(\rho_{n,n}=1\)。其逆极限是
	\[
		\varprojlim_n E_n
		=\left\{(e_n)\in\prod_{n\geq1}E_n:
		\rho_n(e_{n+1})=e_n\text{ 对每个 }n\text{ 成立}\right\}.
	\]
	逆系统之间的同态是一族与转移同态交换的同态；它逐个分量诱导逆极限上的同态。
\end{definition}

逆极限中的元素保存了所有层上的数据，以及这些数据之间的相容性。给定同态 \(f_n:T\to E_n\)，若 \(\rho_n f_{n+1}=f_n\)，就有唯一的同态 \(T\to\varprojlim E_n\)，把 \(t\) 送到 \((f_n(t))_n\)。这个逐项描述也证明了逆极限的泛性质：
\[
	\begin{tikzcd}[column sep=large]
		& T \arrow[dl,"f_{n+1}"'] \arrow[d,dashed,"\exists!"] \arrow[dr,"f_n"] & \\
		E_{n+1} & \varprojlim E_j \arrow[l,"\pi_{n+1}"] \arrow[r,"\pi_n"'] & E_n
	\end{tikzcd}
\]
图中的投影另满足 \(\rho_n\pi_{n+1}=\pi_n\)。环的逆极限也用同样的相容族构造，运算逐项进行。

\begin{definition}\label{g5:def:mittag-leffler}
	设 \((E_n,\rho_n)\) 为可数逆系统。若对每个固定的 \(n\)，下降链
	\[
		E_n\supseteq\operatorname{im}(E_{n+1}\to E_n)
		\supseteq\operatorname{im}(E_{n+2}\to E_n)\supseteq\cdots
	\]
	最终稳定，就称该系统满足\term[MittagLefflertiaojian]{Mittag–Leffler 条件}[Mittag–Leffler condition]，简称 ML 条件。
\end{definition}

这里稳定的阶数可以随 \(n\) 改变。每个转移同态都满射的系统满足 ML 条件；若每个 \(E_n\) 都有有限长度，其子模下降链稳定，也满足这一条件。它控制的是从更高阶能够传下来的像，正好适于调整提升。\footnote{Magnus Gustaf（Gösta）Mittag-Leffler（1846--1927），瑞典数学家，研究解析函数论，并创办 Acta Mathematica。}

\begin{lemma}\label{g5:lem:inverse-limit-exact}
	设 \(A\) 为交换含幺环，且有可数 \(A\) 模逆系统之间的逐项短正合列
	\[
		0\longrightarrow L_n\xrightarrow{\iota_n}E_n
		\xrightarrow{p_n}Q_n\longrightarrow0,
	\]
	其中各转移同态与这两族映射交换。则
	\[
		0\longrightarrow\varprojlim L_n
		\longrightarrow\varprojlim E_n
		\longrightarrow\varprojlim Q_n
	\]
	正合。若核系统 \((L_n)\) 满足 ML 条件，最后一张映射还满射。特别地，核系统的全部转移同态满射时，可对短正合列取逆极限而保持短正合。
\end{lemma}
\begin{proof}
	第一张映射的单射性逐个分量可见。若相容族 \((e_n)\) 的像为零，则各 \(e_n\) 唯一来自 \(l_n\in L_n\)。转移相容性及 \(\iota_n\) 的单射性保证 \((l_n)\) 相容，所以中间正合。

	为证附加的满射性，固定 \(q=(q_n)\in\varprojlim Q_n\)。令
	\[
		T_n=\{e\in E_n:p_n(e)=q_n\}.
	\]
	每个 \(T_n\) 都非空，并且任取其中一点后，整个 \(T_n\) 就是它加上 \(\iota_n(L_n)\) 得到的陪集。转移同态把 \(T_m\) 映入 \(T_n\)。固定 \(n\)，当 \(m\) 增大时，像 \(T_{m,n}\subseteq T_n\) 依次缩小；它是子群 \(\operatorname{im}(L_m\to L_n)\) 的一个非空陪集，这里将 \(L_n\) 与其在 \(E_n\) 中的像等同。

	由 ML 条件，这些子群在某阶以后相同。相同子群的两个非空陪集若一个包含另一个，它们必相等。因此 \(T_{m,n}\) 也最终稳定，记稳定的非空集合为 \(T_n'\)。我们证明 \(T_{n+1}'\to T_n'\) 满射。给定 \(e\in T_n'\)，取充分大的 \(m\)，使从 \(T_m\) 到第 \(n\)、\(n+1\) 层的像都已稳定。因为 \(T_{m,n}=T_n'\)，可选 \(\widetilde e\in T_m\) 映到 \(e\)；它在第 \(n+1\) 层的像属于 \(T_{n+1}'\)，并继续映到 \(e\)。反过来，\(T_{n+1}'\) 中任一点也来自这样一个充分高的 \(T_m\)，故其在第 \(n\) 层的像必落在 \(T_n'\)。

	先取 \(e_1\in T_1'\)，再利用这些满射逐次选取 \(e_{n+1}\in T_{n+1}'\)，使它映到已选的 \(e_n\)。由此得到 \((e_n)\in\varprojlim E_n\)，其像恰为 \(q\)。
\end{proof}

若核系统的转移同态已经满射，上述选择也可以直接理解：下一阶的任意提升与已选提升之差属于当前核，将这个差从上一阶的核提上去，再减掉它即可。ML 条件把同样的调整限制在最终稳定的像中。

\begin{example}\label{g5:exm:inverse-limit-not-right-exact}
	取逐项短正合列
	\[
		0\longrightarrow2^n\mathbb Z\longrightarrow\mathbb Z
		\longrightarrow\mathbb Z/2^n\mathbb Z\longrightarrow0.
	\]
	左侧转移同态为包含，中间为恒等，右侧为模较低次幂的映射。左侧逆极限是 \(\bigcap_n2^n\mathbb Z=0\)，中间逆极限是 \(\mathbb Z\)。一般地，对素数 \(p\)，环
	\[
		\mathbb Z_p=\varprojlim_n\mathbb Z/p^n\mathbb Z
	\]
	称为\term[pjinzhengshuhuan]{\(p\) 进整数环}[ring of p-adic integers]。在本例中，极限列的最后一张映射是 \(\mathbb Z\to\mathbb Z_2\)。

	为看出它不满射，令 \(\alpha_n\in\mathbb Z/2^n\mathbb Z\) 为有限和
	\[
		\sum_{\substack{j\geq0\\2j<n}}2^{2j}
	\]
	的剩余类。\((\alpha_n)\) 相容，确定 \(\alpha\in\mathbb Z_2\)。有限几何级数公式给出 \(3\alpha_n=-1\pmod{2^n}\)，故 \(3\alpha=-1\)。若它来自整数 \(a\)，则整数 \(3a+1\) 被每个 \(2^n\) 整除，因而为零；但不存在满足 \(3a=-1\) 的整数。因此这张极限映射不满射。核系统在第 \(n\) 层的像是持续严格下降的 \(2^m\mathbb Z\)、\(m\geq n\)，正好不满足 ML 条件。
\end{example}

\subsection{完备化的正合性}
\label{g5:subsec:0703:02}

设 \(A\) 为交换含幺环，\(I\subseteq A\) 为理想。对 \(A\) 模 \(M\)，记其 \(I\) 进完备化为
\[
	\widehat M=\varprojlim_{n\geq1}M/I^nM,
	\qquad
	\eta_M:M\longrightarrow\widehat M,
	\quad m\longmapsto(m+I^nM)_n.
\]
本节所有带帽号都取同一个理想 \(I\)。\(\widehat A\) 是环，且 \(\widehat M\) 自然为 \(\widehat A\) 模：在第 \(n\) 层用 \(A/I^n\) 作用于 \(M/I^nM\)，再逐项相乘。\(A\) 线性映射 \(f:M\to M'\) 保持各 \(I^n\) 子模，因而诱导 \(\widehat A\) 线性映射 \(\widehat f\)。有限直和与完备化相容，因为相容条件可逐个分量检查。

\ProseParagraphBreak

投影到第 \(n\) 层的核记为
\[
	F^n\widehat M=\ker(\widehat M\to M/I^nM).
\]
这些子模规定完备化的拓扑。\(\eta_M\) 的核为 \(\bigcap_n I^nM\)，所以原模是否嵌入其完备化取决于原拓扑是否分离。\(F^n\widehat M\) 与 \(I^n\widehat M\) 的相等性将在有限模情形得到，当前暂用投影核描述拓扑。

\begin{definition}\label{g5:def:adic-completeness}
	设 \(A\) 为交换含幺环，\(I\subseteq A\) 为理想，\(M\) 为 \(A\) 模。若序列 \((m_j)_{j\geq1}\) 满足：对每个 \(n\)，存在 \(j_0\)，使
	\[
		j,k\geq j_0\quad\Longrightarrow\quad m_j-m_k\in I^nM,
	\]
	则称它为 \(I\) 进拓扑中的\term[Cauchylie]{Cauchy 列}[Cauchy sequence]。\footnote{Augustin-Louis Cauchy（1789--1857），法国数学家，在分析的严格化和复函数理论方面作出重要贡献。}若自然映射 \(\eta_M:M\to\widehat M\) 是同构，则称 \(M\) 对 \(I\) 进拓扑完备且分离。
\end{definition}

后一条件等价于拓扑分离且每个 Cauchy 列都有极限。事实上，Cauchy 列在每个 \(M/I^nM\) 中最终稳定；这些稳定值相容，若 \(\eta_M\) 为同构，就确定唯一极限。反过来，给定相容族 \((\overline m_n)\)，任取代表 \(m_n\in M\)，便有 \(m_{n+1}-m_n\in I^nM\)，因而形成 Cauchy 列；它的极限在每个商中等于指定的 \(\overline m_n\)，给出 \(\eta_M\) 的满射性，而分离性给出单射性。同一解释适用于由任意子模降链规定的拓扑。

\begin{theorem}\label{g5:thm:completion-exact}
	设 \(A\) 为 Noether 环，\(I\subseteq A\) 为任意理想。对有限生成 \(A\) 模的短正合列
	\[
		0\longrightarrow N\longrightarrow M\longrightarrow Q\longrightarrow0,
	\]
	取 \(I\) 进完备化后得到短正合列
	\[
		0\longrightarrow\widehat N\longrightarrow\widehat M
		\longrightarrow\widehat Q\longrightarrow0.
	\]
\end{theorem}
\begin{proof}
	将 \(N\) 看作 \(M\) 的子模。对每个 \(n\geq1\)，真正保持左端单射的商列是
	\[
		0\longrightarrow N/(N\cap I^nM)
		\longrightarrow M/I^nM
		\longrightarrow Q/I^nQ\longrightarrow0.
	\]
	其中最后一个商等于 \(M/(N+I^nM)\)。左侧核系统的转移同态全部满射，所以\cref{g5:lem:inverse-limit-exact}给出
	\[
		0\longrightarrow\varprojlim_n N/(N\cap I^nM)
		\longrightarrow\widehat M\longrightarrow\widehat Q\longrightarrow0.
	\]

	剩下的任务是将左端识别为 \(\widehat N\)。由 Artin–Rees 的拓扑推论 \cref{g5:cor:induced-adic-topology}，存在 \(c\geq0\)，使对每个 \(n\geq1\) 有
	\[
		N\cap I^{n+c}M\subseteq I^nN\subseteq N\cap I^nM.
	\]
	右侧包含诱导映射
	\(\widehat N\to\varprojlim N/(N\cap I^nM)\)。反向映射在第 \(n\) 层，取相容族的第 \(n+c\) 项，再按左侧包含映到 \(N/I^nN\)。改变代表元不影响结果，相邻层的相容性来自原族的相容性。

	两张映射复合后，在第 \(n\) 层都是把第 \(n+c\) 项降到第 \(n\) 层；这就是原族的第 \(n\) 项。因此它们互逆，且所得识别与 \(N\to M\) 相容。代入前面的极限短正合列，便得到结论。
\end{proof}

这里 Artin–Rees 解决的是左端所用拓扑的问题：\(N/I^nN\to M/I^nM\) 可能不单射，但诱导滤过 \(N\cap I^nM\) 与本来的 \(I^nN\) 互相共尾，从而具有同一个完备化。可数核系统的提升引理则负责右端满射。这两项作用在证明中分别出现；对应结果可对照\cite[第8节，定理8.1、8.6，第56—60页]{Matsumura1989}。

\subsection{有限模的张量表示}
\label{g5:subsec:0703:03}

\(\widehat M\) 已有 \(\widehat A\) 模结构，因而存在自然映射
\[
	\theta_M:\widehat A\otimes_A M\longrightarrow\widehat M,
	\qquad a\otimes m\longmapsto a\eta_M(m).
\]
右侧的表达式对 \(A\) 平衡，故确实定义了张量积上的同态。有限呈示将它的同构性归结为有限自由模的情形。

\begin{theorem}\label{g5:thm:completion-tensor}
	设 \(A\) 为 Noether 环，\(I\subseteq A\) 为理想，\(M\) 为有限生成 \(A\) 模。则自然映射
	\[
		\theta_M:\widehat A\otimes_A M\xrightarrow{\ \cong\ }\widehat M
	\]
	是 \(\widehat A\) 模同构，并与 \(A\) 线性映射相容。特别地，\(\widehat M\) 是有限生成 \(\widehat A\) 模。
\end{theorem}
\begin{proof}
	由\cref{g5:prop:noetherian-finite-modules}，\(M\) 有有限呈示
	\[
		A^u\xrightarrow{D}A^v\longrightarrow M\longrightarrow0.
	\]
	把这列在 \(\operatorname{im}D\) 处分成两条短正合列；其中涉及的核与像仍有限生成。由\cref{g5:thm:completion-exact}，其完备化列右正合。另一方面，\cref{g5:prop:tensor-right-exact}给出张量积行的右正合性，得到交换图
	\[
		\begin{tikzcd}[column sep=large]
			\widehat A\otimes_A A^u \arrow[r] \arrow[d,"\cong"'] &
			\widehat A\otimes_A A^v \arrow[r] \arrow[d,"\cong"'] &
			\widehat A\otimes_A M \arrow[r] \arrow[d,"\theta_M"] & 0 \\
			\widehat{A^u} \arrow[r] & \widehat{A^v} \arrow[r] & \widehat M \arrow[r] & 0.
		\end{tikzcd}
	\]
	有限直和与完备化相容，使左侧两张竖箭头分别成为 \(\widehat A^u\)、\(\widehat A^v\) 的自然识别。两行的第一张箭头也都由矩阵 \(D\) 在 \(\widehat A\) 中的像给出，所以最右侧两个模是同一映射的余核，\(\theta_M\) 必为同构。

	若 \(f:M\to M'\) 为 \(A\) 线性映射，则在纯张量上有
	\(\widehat f(a\eta_M(m))=a\eta_{M'}(f(m))\)，证明自然性。有限个 \(M\) 生成元在 \(\theta_M\) 下的像生成 \(\widehat M\)，得到最后的有限性。
\end{proof}

这项表示把“所有相容截断”变成一次基环扩张，见\cite[第8节，定理8.7，第60页]{Matsumura1989}。若 \(A\) 本身 \(I\) 进完备，即 \(A\to\widehat A\) 是同构，则每个有限生成 \(A\) 模也完备，因为上述自然同构此时就是 \(M\to\widehat M\)。

\begin{example}\label{g5:exm:completion-infinite-direct-sum}
	设 \(k\) 为域，\(A=k[t]\)，\(I=(t)\)，并令
	\[
		M=\bigoplus_{j\geq0}Ae_j.
	\]
	此时 \(\widehat A=k[[t]]\)。对每个 \(n\)，有限和
	\[
		x_n=\sum_{j=0}^{n-1}t^je_j\pmod{t^nM}
	\]
	形成相容族，故确定 \(x\in\widehat M\)。但
	\[
		\widehat A\otimes_A M\cong\bigoplus_{j\geq0}k[[t]]e_j
	\]
	中的元素只有有限多个非零坐标。将 \(\widehat M\) 投影到第 \(j\) 个坐标的完备化，\(x\) 的像为非零级数 \(t^j\)。因此 \(x\) 不在 \(\theta_M\) 的像中。这说明对无限生成模，张量表示可能失效；无限直和也不能直接与完备化交换。
\end{example}

\subsection{完备化环的 Noether 性}
\label{g5:subsec:0703:04}

先确定 \(\widehat A\) 上的理想幂究竟给出哪些截断。这一步也把逆极限的拓扑识别为通常的理想进拓扑，供随后逐阶逼近时使用。

\begin{proposition}\label{g5:prop:completion-quotient}
	设 \(A\) 为 Noether 环，\(I\subseteq A\) 为理想，\(M\) 为有限生成 \(A\) 模，所有完备化均取 \(I\) 进完备化。对任意理想 \(H\subseteq A\)，有自然同构
	\[
		\widehat{M/HM}\cong\widehat M/H\widehat M.
	\]
	特别地，对每个 \(n\geq1\)，自然投影诱导
	\[
		\widehat M/I^n\widehat M\cong M/I^nM,
		\qquad
		\widehat A/I^n\widehat A\cong A/I^n.
	\]
	令 \(J=I\widehat A\)。则 \(F^n\widehat M=J^n\widehat M\)，并且 \(\widehat M\) 对 \(J\) 进拓扑完备且分离。
\end{proposition}
\begin{proof}
	对 \(0\to HM\to M\to M/HM\to0\) 使用完备化正合性。通过\cref{g5:thm:completion-tensor}及其自然性，\(\widehat{HM}\to\widehat M\) 的像等于
	\[
		\widehat A\cdot\eta_M(HM)=H\widehat M.
	\]
	具体地，\(HM\) 的元素是有限和 \(\sum_i h_im_i\)，其任意 \(\widehat A\) 倍数均落在 \(H\widehat M\) 中；反过来，\(\widehat M\) 由所有 \(\eta_M(m)\) 在 \(\widehat A\) 上生成，所以 \(H\widehat M\) 的每个元素也在该像中。取余核就得到第一个同构。

	令 \(H=I^n\)。模 \(M/I^nM\) 被 \(I^n\) 零化，其 \(I\) 进商系统从第 \(n\) 层以后为常值系统，故其完备化就是它自身。于是第一个同构化为
	\(\widehat M/I^n\widehat M\cong M/I^nM\)，而相应映射正是第 \(n\) 个投影。又因
	\(J^n\widehat M=I^n\widehat M\)，这证明了投影核的描述。

	一个相容族若每个投影都为零，它本身就为零，因此
	\[
		\bigcap_{n\geq1}J^n\widehat M=0.
	\]
	上面的商同构与各转移映射相容，故自然映射
	\[
		\widehat M\longrightarrow
		\varprojlim_n\widehat M/J^n\widehat M
		\cong\varprojlim_n M/I^nM=\widehat M
	\]
	是恒等识别。这同时证明完备性及分离性，也保证每个 \(J\) 进 Cauchy 列收敛到唯一极限。取 \(M=A\) 就得到环的全部结论。
\end{proof}

这些商同构及拓扑识别对应于\cite[第8节，定理8.11、8.13，第61页]{Matsumura1989}。其中并未使用 \(\widehat A\) 的 Noether 性；它们说明有限阶方程及理想幂在完备化后保持原有形式。

\begin{proposition}\label{g5:prop:completion-associated-graded}
	设 \(A\) 为 Noether 环，\(I\subseteq A\) 为理想，并令 \(J=I\widehat A\)。自然同态 \(A\to\widehat A\) 诱导分次环同构
	\[
		\operatorname{gr}_I A
		=\bigoplus_{n\geq0}I^n/I^{n+1}
		\xrightarrow{\ \cong\ }
		\bigoplus_{n\geq0}J^n/J^{n+1}
		=\operatorname{gr}_J\widehat A.
	\]
\end{proposition}
\begin{proof}
	次数零的同构为 \(A/I\cong\widehat A/J\)。对 \(n\geq1\)，\cref{g5:prop:completion-quotient}给出与自然满射相容的两层同构
	\[
		\begin{tikzcd}
			A/I^{n+1} \arrow[r] \arrow[d,"\cong"'] & A/I^n \arrow[d,"\cong"] \\
			\widehat A/J^{n+1} \arrow[r] & \widehat A/J^n.
		\end{tikzcd}
	\]
	取水平映射的核，就得到 \(I^n/I^{n+1}\cong J^n/J^{n+1}\)。它们全由同一环同态诱导，故与各次数间的乘法相容。
\end{proof}

\begin{theorem}\label{g5:thm:completion-noetherian}
	设 \(A\) 为 Noether 环，\(I\subseteq A\) 为任意理想。则其 \(I\) 进完备化 \(\widehat A\) 也是 Noether 环。
\end{theorem}
\begin{proof}
	记 \(R=\widehat A\)、\(J=IR\)。由\cref{g5:prop:completion-quotient}，\(R\) 对 \(J\) 进拓扑完备且分离。由\cref{g5:prop:rees-associated-graded,g5:prop:completion-associated-graded}，\(\operatorname{gr}_J R\) 为 Noether 环：若 \(I=(a_1,\ldots,a_r)\)，则 \(\operatorname{gr}_I A\) 由这些元素的次数一初始类在 Noether 环 \(A/I\) 上生成，是有限变量多项式环的商。

	固定 \(R\) 的任意理想 \(\mathfrak a\)。我们先为它选择有限个元素，使其初始类生成全部初始关系，再把逐阶消去所得的系数收敛起来。定义分次理想
	\[
		\mathfrak a^*
		=\bigoplus_{d\geq0}
		\frac{(\mathfrak a\cap J^d)+J^{d+1}}{J^{d+1}}
		\subseteq\operatorname{gr}_J R.
	\]
	它在次数 \(d\) 的元素，正是 \(\mathfrak a\cap J^d\) 中元素的初始类。乘以次数 \(e\) 的元素后，仍来自 \(\mathfrak a\cap J^{d+e}\)，所以它确为分次理想。

	由于 \(\operatorname{gr}_J R\) Noether，\(\mathfrak a^*\) 有有限个齐次生成元。这里可以先取任意有限生成元，再取其全部非零齐次分量；分次理想包含这些分量，因此仍得到有限生成元。选取对应的
	\[
		f_1,\ldots,f_q\in\mathfrak a,
		\qquad
		f_i\in J^{d_i}\setminus J^{d_i+1},
	\]
	使它们的初始类生成 \(\mathfrak a^*\)。若 \(\mathfrak a^*=0\)，分离性使 \(\mathfrak a=0\)，此时结论成立。以下设 \(q>0\)。

	任取非零 \(f\in\mathfrak a\)。因为 \(\bigcap_dJ^d=0\)，存在有限的 \(d\)，使 \(f\in J^d\setminus J^{d+1}\)。在分次环的次数 \(d\) 部分，可以将其初始类写成
	\[
		f^*=\sum_{d_i\leq d} b_i^* f_i^*,
		\qquad
		b_i^*\in J^{d-d_i}/J^{d-d_i+1}.
	\]
	这是把一个生成元表示式取次数 \(d\) 的齐次部分。选择提升 \(b_i\in J^{d-d_i}\)，并对 \(d_i>d\) 取 \(b_i=0\)，就有
	\[
		f-\sum_i b_i f_i\in\mathfrak a\cap J^{d+1}.
	\]
	若余项非零，就继续对它作同样的消去。

	若过程在有限步结束，已得 \(f\in(f_1,\ldots,f_q)\)。否则，设各步非零余项的阶为严格递增的 \(d_0,d_1,\ldots\)，第 \(\nu\) 步的修正系数为 \(b_{i,\nu}\)。当 \(d_\nu\geq d_i\) 时，
	\[
		b_{i,\nu}\in J^{d_\nu-d_i}.
	\]
	因此对每个 \(i\)，部分和 \(\sum_{\nu<k}b_{i,\nu}\) 是 \(J\) 进 Cauchy 列：任意指定的 \(J^N\)，最终都包含每一个后续修正系数及其有限和。由 \(R\) 的完备性，它收敛到某个 \(b_i\in R\)。逐次消去还给出
	\[
		f-\sum_i\left(\sum_{\nu<k}b_{i,\nu}\right)f_i
		\in J^{d_k}.
	\]
	对任意 \(N\)，取充分大的 \(k\)，两边按模 \(J^N\) 稳定，得到 \(f-\sum_i b_i f_i\in J^N\)。分离性于是给出
	\[
		f=\sum_{i=1}^q b_i f_i.
	\]
	最终表达式仍只用有限个 \(f_i\)，故 \(\mathfrak a=(f_1,\ldots,f_q)\)。任意理想均有限生成，证明 \(R\) 为 Noether 环。
\end{proof}

上述结论亦见\cite[第8节，定理8.12，第61页]{Matsumura1989}。本证明中的收敛发生在有限组系数上：初始理想的有限生成控制每一阶需要使用的元素，而完备性允许把所有阶的修正合成为同一组系数。由此，形式级数所容纳的无限展开与 Noether 性所要求的有限生成并不冲突。
% 纲要条目排为编号小节标题，并进入目录与书签；说明段仍保留于纲要注释中。

% !TeX root = ../../Book5.tex
% 纲要标题：### 7.4 忠实平坦及局部结构
% 纲要位置：几何学计划纲要.md，第 5283 行。
% 正文状态：初稿；尚未正式审核。

% 纲要内容开始
% 
% * 理想张量判据推出平坦
% * 剩余域不变与忠实性
% * 维数和伴随分次比较
% * 扩张收缩与局部检测
% * 完备化不是局部化或 Hensel 化
% 
% 纲要内容结束

\section{忠实平坦及局部结构}
\label{g5:sec:0704}

前一节已经证明，有限模的完备化既保持正合列，又等于与完备化环张量。把这两件事放在一起，便得到完备化的平坦性。对于局部环，这个扩张还保留剩余域，从而忠实平坦。我们将用这些事实比较理想、维数和一阶结构，再说明补齐无穷级数与反演分母、提升代数方程的根各自完成什么工作。

\subsection{理想张量判据推出平坦}
\label{g5:subsec:0704:01}

\begin{theorem}\label{g5:thm:completion-flat}
	设 \(A\) 为 Noether 环，\(I\subseteq A\) 为理想，\(\widehat A=\varprojlim_{n\geq1}A/I^n\)。则自然同态 \(A\to\widehat A\) 平坦。
\end{theorem}
\begin{proof}
	由\cref{g5:thm:flat-ideal-criterion}，只需证明任意有限生成理想 \(J\subseteq A\) 的包含张量后仍单射。Noether 性保证 \(J\) 和 \(A/J\) 都为有限模。对
	\[
		0\longrightarrow J\longrightarrow A\longrightarrow A/J\longrightarrow0
	\]
	应用\cref{g5:thm:completion-exact}，得到单射 \(\widehat J\to\widehat A\)。再由自然的张量表示\cref{g5:thm:completion-tensor}，这张单射正是
	\[
		J\otimes_A\widehat A\longrightarrow A\otimes_A\widehat A=\widehat A.
	\]
	因此理想判据的全部要求成立。
\end{proof}

这条证明的有限性落在 \(J\) 上：完备化保持的是有限模正合列，而平坦性要求张量保持任意模单射。第六章的理想判据把两种范围接在了一起。结果及同一证明思路见\cite[第8节，定理8.8，第60页]{Matsumura1989}。

\ProseParagraphBreak

平坦性本身并不保证 \(A\to\widehat A\) 单射。例如取 \(A=k\times k\)、\(I=k\times0\)，有 \(I^n=I\)，所以 \(\widehat A=A/I\cong k\)。完备化在这里是一个平坦投影。下一步需要确定，何时它还保留对全部模的检测能力。

\subsection{剩余域不变与忠实性}
\label{g5:subsec:0704:02}

\begin{theorem}\label{g5:thm:completion-faithful}
	设 \(A\) 为 Noether 环，\(I\subseteq A\) 为理想。则 \(A\to\widehat A\) 忠实平坦，当且仅当
	\[
		I\subseteq\operatorname{Jac}(A).
	\]
	特别地，若 \((A,\mathfrak m,\kappa)\) 为非零 Noether 局部环，且完备化取 \(\mathfrak m\) 进拓扑，则 \(\widehat A\) 为 Noether 局部环，其极大理想和剩余域分别为
	\[
		\widehat{\mathfrak m}=\mathfrak m\widehat A,
		\qquad \widehat A/\widehat{\mathfrak m}\cong\kappa,
	\]
	且 \(A\to\widehat A\) 忠实平坦。
\end{theorem}
\begin{proof}
	对任意极大理想 \(\mathfrak p\subset A\)，有限模张量表示给出
	\[
		\widehat A/\mathfrak p\widehat A
		\cong\widehat{A/\mathfrak p}.
	\]
	若 \(I\subseteq\mathfrak p\)，右侧的 \(I\) 作用为零，其完备化就是非零域 \(A/\mathfrak p\)。若 \(I\nsubseteq\mathfrak p\)，\(I\) 在这个域中的像为单位理想，全部截断都是零，右侧因而为零。由上一条平坦性定理及\cref{g5:thm:faithful-flat-spectrum}，忠实性恰好要求第一种情形对每个极大理想成立，即 \(I\subseteq\operatorname{Jac}(A)\)。

	现在令 \(I=\mathfrak m\)。\cref{g5:thm:completion-noetherian}给出 \(\widehat A\) Noether；\cref{g5:prop:completion-quotient}给出
	\(\widehat A/\mathfrak m\widehat A\cong A/\mathfrak m=\kappa\)，所以 \(\mathfrak m\widehat A\) 是极大理想。若 \(b\notin\mathfrak m\widehat A\)，选 \(a\in A\) 使它在这个商域中与 \(b\) 同余。此时 \(a\notin\mathfrak m\)，故 \(a\) 在 \(A\) 中可逆，并可写
	\[
		b=a(1-u),\qquad u\in\mathfrak m\widehat A.
	\]
	由完备性，\(\sum_{j\geq0}u^j\) 收敛；部分和与 \(1-u\) 的乘积为 \(1-u^{N+1}\)，极限为 \(1\)。因此 \(b\) 可逆，说明所给极大理想唯一。最后 \(\mathfrak m=\operatorname{Jac}(A)\)，第一部分给出忠实平坦性。
\end{proof}

一般理想进完备化的忠实性判据及局部总结可对照\cite[第8节，定理8.14及节末总结，第62—63页]{Matsumura1989}。局部情形中，自然同态因而单射；这也与 Krull 交定理给出的 \(\bigcap_n\mathfrak m^n=0\) 一致。

\subsection{维数和伴随分次比较}
\label{g5:subsec:0704:03}

完备化增加了可以同时拥有无穷多个系数的元素，却没有改变任何固定阶的截断。因此，从有限阶截断读取的伴随分次环、嵌入维数和长度函数应当保持不变。Krull 维数由素链定义，需要另外给出比较依据。

\begin{proposition}\label{g5:prop:completion-graded-comparison}
	设 \((A,\mathfrak m,\kappa)\) 为非零 Noether 局部环，\(M\) 为有限 \(A\) 模。取极大理想进完备化，令 \(\widehat{\mathfrak m}=\mathfrak m\widehat A\)。则存在自然同构
	\[
		\operatorname{gr}_{\mathfrak m}A
		\cong\operatorname{gr}_{\widehat{\mathfrak m}}\widehat A,
		\qquad
		\operatorname{gr}_{\mathfrak m}M
		\cong\operatorname{gr}_{\widehat{\mathfrak m}}\widehat M.
	\]
	这些同构保留分次及环、模的乘法。特别地，
	\[
		\operatorname{edim}A=\operatorname{edim}\widehat A,\qquad
		\ell_A(M/\mathfrak m^{n+1}M)
		=\ell_{\widehat A}(\widehat M/\widehat{\mathfrak m}^{\,n+1}\widehat M)
		\quad(n\geq0).
	\]
\end{proposition}
\begin{proof}
	环的分次同构已由\cref{g5:prop:completion-associated-graded}证明。对模使用同样的相邻层核比较：\cref{g5:prop:completion-quotient}的商同构与截断映射相容，比较
	\[
		M/\mathfrak m^{n+1}M\longrightarrow M/\mathfrak m^nM
	\]
	和完备化后的同一映射，并取核，得到
	\[
		\mathfrak m^nM/\mathfrak m^{n+1}M
		\cong
		\widehat{\mathfrak m}^{\,n}\widehat M/
		\widehat{\mathfrak m}^{\,n+1}\widehat M.
	\]
	当 \(M=A\) 时，同构由自然环同态诱导，因此保持齐次乘法；对一般 \(M\) 也保持环对模的作用。取 \(n=1\)、\(M=A\)，并使用两环相同的剩余域，就得到嵌入维数相等。最后，长度可按有限滤过的加性计算，每个分次商被极大理想零化，其长度正是作为 \(\kappa\) 向量空间的维数，故两侧长度相等。
\end{proof}

\begin{lemma}\label{g5:lem:flat-going-down}
	设 \(A\to B\) 为平坦的交换含幺环同态。若
	\(\mathfrak p\subseteq\mathfrak p'\) 为 \(A\) 的素理想，
	\(\mathfrak q'\) 为 \(B\) 的素理想且收缩为 \(\mathfrak p'\)，则存在
	\(\mathfrak q\subseteq\mathfrak q'\)，其收缩为 \(\mathfrak p\)。
\end{lemma}
\begin{proof}
	相对局部检测\cref{g5:prop:flat-local-test}说明
	\(A_{\mathfrak p'}\to B_{\mathfrak q'}\) 平坦。它是非零局部环之间的局部同态，由\cref{g5:cor:flat-local-faithful}忠实平坦，故谱映射满射。于是 \(A_{\mathfrak p'}\) 的素理想 \(\mathfrak pA_{\mathfrak p'}\) 有一个上方素理想。局部化的素理想对应把它送回所需的 \(\mathfrak q\subseteq\mathfrak q'\)，收缩再取局部化对应便为 \(\mathfrak p\)。
\end{proof}

\begin{theorem}\label{g5:thm:completion-dimension}
	设 \((A,\mathfrak m)\) 为非零 Noether 局部环，\(\widehat A\) 为其极大理想进完备化。则
	\[
		\dim\widehat A=\dim A.
	\]
	因此 \(A\) 为正则局部环，当且仅当 \(\widehat A\) 为正则局部环。
\end{theorem}
\begin{proof}
	设 \(d=\dim A\)。由\cref{g5:thm:system-of-parameters}，存在 \(d\) 个元素 \(a_1,\ldots,a_d\in\mathfrak m\)，其生成理想 \(J\) 满足 \(\sqrt J=\mathfrak m\)。因为 \(\mathfrak m\) 有限生成，存在 \(N\) 使 \(\mathfrak m^N\subseteq J\)：给每个生成元选一个落入 \(J\) 的幂，再对足够高次单项式使用抽屉原理。扩张到 \(\widehat A\) 得
	\[
		\widehat{\mathfrak m}^{\,N}
		\subseteq J\widehat A\subseteq\widehat{\mathfrak m}.
	\]
	所以 \(J\widehat A\) 的根式为 \(\widehat{\mathfrak m}\)，它仍由 \(d\) 个元素生成。对 Noether 环 \(\widehat A\) 应用\cref{g5:thm:height-theorem}，得到 \(\dim\widehat A\leq d\)。\(d=0\) 时相同论证采用空生成元表。

	反方向，取 \(A\) 中一条长度为 \(d\) 的素链
	\[
		\mathfrak p_0\subsetneq\cdots\subsetneq\mathfrak p_d=\mathfrak m.
	\]
	维数有限，故这样的链存在，且最长链的末项必为极大理想。以收缩为 \(\mathfrak m\) 的 \(\widehat{\mathfrak m}\) 为末端，逐次应用上个引理，得到 \(\widehat A\) 中收缩为整条给定素链的素链。相邻收缩不同，提升后的包含也严格，所以 \(\dim\widehat A\geq d\)。

	最后，前一命题已经证明嵌入维数不变。两环都为 Noether 局部环，因此定义中的等式 \(\dim A=\operatorname{edim}A\) 在完备化前后等价。
\end{proof}

这里对维数的比较用了参数系和素链提升；长度函数的相同则来自逐阶商同构。二者共同说明，切锥的分次数据与局部维数都保留了下来。局部分支却可能增加：第五章的\cref{g5:exm:formal-node}给出一个整的结点局部环，其完备化同构于 \(k[[u,v]]/(uv)\)，具有两个极小素理想。因此维数不变并不要求素谱上的点一一对应。

\subsection{扩张收缩与局部检测}
\label{g5:subsec:0704:04}

\begin{proposition}\label{g5:prop:completion-submodule-contraction}
	设 \(A\) 为 Noether 环，\(I\subseteq\operatorname{Jac}(A)\)，\(M\) 为有限 \(A\) 模，\(N\subseteq M\)。把 \(M\) 嵌入 \(\widehat M\)，并把 \(\widehat N\) 识别为其中的子模，则
	\[
		M\cap\widehat N=N.
	\]
	特别地，对任意理想 \(J\subseteq A\)，有 \(J\widehat A\cap A=J\)。任意 \(A\) 模序列的正合性，以及同态的单射、满射和同构性质，都可在张量 \(\widehat A\) 后检测。
\end{proposition}
\begin{proof}
	\cref{g5:thm:completion-faithful}及\cref{g5:thm:faithful-flat-equalizer}保证任意 \(A\) 模 \(L\) 的自然映射
	\(L\to\widehat A\otimes_A L\) 单射。将此用于 \(M\) 和 \(M/N\)，再用完备化正合列
	\[
		0\longrightarrow\widehat N\longrightarrow\widehat M
		\longrightarrow\widehat{M/N}\longrightarrow0,
	\]
	可知 \(m\in M\) 的像属于 \(\widehat N\)，恰好等价于 \(m+N\) 在 \(\widehat{M/N}\) 中为零，也就是 \(m\in N\)。取 \(M=A,N=J\)，张量表示把 \(\widehat J\) 的像识别为 \(J\widehat A\)。最后一项是\cref{g5:prop:faithful-flat-reflects-exactness}的应用；它对任意模成立，无需有限生成。
\end{proof}

注意最后一句的操作是与 \(\widehat A\) 张量。对无限生成模，尚不能用本章的有限模张量表示，把它随意改写为该模的逆极限完备化。下面的有限自由检测则完全落在已证范围内。

\begin{corollary}\label{g5:cor:completion-free-detection}
	设 \((A,\mathfrak m,\kappa)\) 为非零 Noether 局部环，\(M\) 为有限 \(A\) 模，\(r\geq0\)。则
	\[
		M\cong A^r\quad\Longleftrightarrow\quad
		\widehat M\cong\widehat A^r.
	\]
\end{corollary}
\begin{proof}
	正方向由有限直和与完备化相容得到。反方向中，剩余域商同构说明
	\(\dim_\kappa M/\mathfrak mM=r\)。提升一组基，Nakayama 引理给出满射 \(A^r\to M\)。完备化后，在选定 \(\widehat M\cong\widehat A^r\) 的识别下，得到一个方阵，其剩余域矩阵可逆，行列式因而是局部环 \(\widehat A\) 的单位。伴随矩阵公式给出该矩阵可逆，所以完备化后的映射为同构。由忠实平坦检测，原映射也为同构。\(r=0\) 时直接使用 Nakayama 引理。
\end{proof}

\subsection{完备化不是局部化或 Hensel 化}
\label{g5:subsec:0704:05}

在 \(A=k[x]_{(x)}\) 中，我们已经反演了所有在原点不为零的多项式；这些分母处理的是原点附近的正则函数。再作完备化得到 \(k[[x]]\)，其中的元素由所有相容截断决定。二者的区别可以不借助大小比较直接看出：任何局部化 \(A\to S^{-1}A\) 若仍使极大理想 \((x)\) 保持为真理想，则 \(S\) 必须避开 \((x)\)；这些元素在 \(A\) 中已经可逆，所以该局部化只给出 \(A\) 自身。因而只要 \(A\ne\widehat A\)，完备化就不可能是这种局部化。

\ProseParagraphBreak

提升单根是完备性的另一项效果。这通常称为 Hensel 引理。\footnote{Kurt Hensel（1861--1941），德国数学家，创立了 \(p\) 进数理论，对局部数论的发展有重要影响。}下面证明其单根版本；更一般的互素因式分解版本见\cite[第8节，定理8.3，第58页]{Matsumura1989}。

\begin{theorem}[Hensel 单根提升]\label{g5:thm:complete-simple-root}
	设 \((R,\mathfrak n)\) 为对 \(\mathfrak n\) 进拓扑完备且分离的局部环，\(F(Y)\in R[Y]\)，\(a_0\in R\)。若
	\[
		F(a_0)\in\mathfrak n,\qquad F'(a_0)\in R^\times,
	\]
	则存在唯一的 \(a\in a_0+\mathfrak n\) 使 \(F(a)=0\)。令
	\[
		a_{j+1}=a_j-F(a_j)F'(a_j)^{-1},
	\]
	则迭代有定义，且
	\(F(a_j)\in\mathfrak n^{2^j}\)、\(a-a_j\in\mathfrak n^{2^j}\)。
\end{theorem}
\begin{proof}
	多项式展开给出
	\[
		F(a+h)=F(a)+F'(a)h+h^2G(a,h)
	\]
	及 \(F'(a+h)-F'(a)\in(h)\)。若 \(a_j\equiv a_0\pmod{\mathfrak n}\)，则 \(F'(a_j)\) 模 \(\mathfrak n\) 与单位 \(F'(a_0)\) 相同，因而可逆。假设
	\(F(a_j)\in\mathfrak n^{2^j}\)，把
	\(h=-F(a_j)F'(a_j)^{-1}\) 代入展开式，一次项恰好消去，得
	\[
		a_{j+1}-a_j\in\mathfrak n^{2^j},\qquad
		F(a_{j+1})\in\mathfrak n^{2^{j+1}}.
	\]
	归纳保证全部迭代有定义，且构成 Cauchy 列。完备性给出极限 \(a\)，多项式运算连续且环分离，故 \(F(a)=0\)。从固定的第 \(j\) 步起，后续差都在 \(\mathfrak n^{2^j}\) 中，逐阶截断说明 \(a-a_j\) 也属于该理想。

	若 \(a,b\in a_0+\mathfrak n\) 都为根，因式分解给出
	\[
		0=F(b)-F(a)=(b-a)\bigl(F'(a)+(b-a)H\bigr).
	\]
	括号内模 \(\mathfrak n\) 为 \(F'(a_0)\) 的非零剩余类，故为单位，得到 \(a=b\)。
\end{proof}

采用单根表述，一个局部环若对每个首一多项式及其剩余域中的单根都能进行上述提升，就称为\term[Henseljubuhuan]{Hensel 局部环}[Henselian local ring]。因此完备局部环是 Hensel 局部环。对 \(A=\mathbb Q[x]_{(x)}\)，多项式 \(Y^2-(1+x)\) 在剩余域中有单根 \(1\)，但在 \(A\) 中没有根：一个有理函数的平方在不可约因子 \(x+1\) 上的阶必为偶数，而 \(1+x\) 的阶为一。它却在 \(\mathbb Q[[x]]\) 中有唯一常数项为 \(1\) 的根。这同时证明 \(A\ne\widehat A\)。

\ProseParagraphBreak

第15章的\subsecref{g5:subsec:1504:03}将构造 Hensel 化 \(A^h\)：它对从 \(A\) 到 Hensel 局部环的局部同态具有初始泛性质，即每张这样的映射都唯一经过 \(A^h\)。这个泛性质与补齐全部 Cauchy 列不同。下面直接证明，在上述有理系数例子中，完备化不具有这一初始性质。

\begin{example}\label{g5:exm:henselian-not-complete}
	令
	\[
		S=\{f\in\mathbb Q[[x]]:f\text{ 在 }\mathbb Q(x)\text{ 上代数}\}.
	\]
	代数元素的和、积及非零元素的倒数仍代数；若 \(f\in S\) 的常数项为零，则 \(f/x\) 仍属于 \(S\)。因此 \(S\) 是包含 \(\mathbb Q[x]_{(x)}\) 的局部环，极大理想为 \(xS\)，剩余域为 \(\mathbb Q\)。

	给定 \(S\) 上首一多项式和剩余域中的单根，在 \(\mathbb Q[[x]]\) 中用上个定理提升，得到根 \(g\)。多项式的有限多个系数生成 \(\mathbb Q(x)\) 的有限扩域 \(L\)，而 \(g\) 在 \(L\) 上代数。代数性的传递性使 \(g\) 在 \(\mathbb Q(x)\) 上代数，故 \(g\in S\)。所以 \(S\) 是 Hensel 局部环。

	\(\mathbb Q(x)\) 上的多项式只有可数多个，每个非零多项式在域 \(\mathbb Q((x))\) 中只有有限多个根，故 \(S\) 可数。另一方面，取系数为 \(0,1\) 的级数已经给出不可数多个 \(\mathbb Q[[x]]\) 元素。因此 \(S\ne\mathbb Q[[x]]\)。它的截断却满足
	\[
		S/x^nS\cong\mathbb Q[x]/(x^n),
	\]
	因为截断由多项式表示，阶至少 \(n\) 的代数级数可除以 \(x^n\) 后仍留在 \(S\)。所以 \(S\) 的完备化正是 \(\mathbb Q[[x]]\)，而 \(S\) 本身不完备。

	若 \(\mathbb Q[[x]]\) 具有上述 Hensel 化泛性质，映射 \(A\to S\) 就会延伸为一个 \(A\) 代数同态 \(\mathbb Q[[x]]\to S\)。这张同态必单射：任一非零级数都写成 \(x^n u\)，其中 \(u\) 为单位，而 \(x\) 在目标中非零、目标为整环，故像不能为零。这与源不可数、目标可数矛盾。因此完备化与 Hensel 化在这个例子中确实不同。
\end{example}

这个例子保留了全部有限阶邻域，又只加入代数的形式元素。下一节将继续比较这些元素与任意形式级数，并把单根提升中的逐阶计算写成方程解的近似过程。
% 纲要条目排为编号小节标题，并进入目录与书签；说明段仍保留于纲要注释中。

% !TeX root = ../../Book5.tex
% 纲要标题：### 7.5 形式方程与有效近似
% 纲要位置：几何学计划纲要.md，第 5291 行。
% 正文状态：初稿；尚未正式审核。

% 纲要内容开始
% 
% * 有限阶邻域和相容族
% * 多项式、代数幂级数与任意形式级数
% * 收敛、代数化及逼近的不同问题
% 
% ---
% 
% 纲要内容结束

\section{形式方程与有效近似}
\label{g5:sec:0705}

一个形式解由无穷多个系数组成，但每次代入方程、比较一个固定阶数时，只使用其中有限多个系数。完备化将这些有限计算组织成相容系统。本节先说明这种相容性如何恢复真正的形式解，再比较多项式、代数级数和一般形式级数。对于 Jacobian 可逆的方程，逐次修正还给出达到指定精度的明确步骤。

\subsection{有限阶邻域和相容族}
\label{g5:subsec:0705:01}

设 \((A,\mathfrak m)\) 为 Noether 局部环。以
\[
	A_n=A/\mathfrak m^n,\qquad n\geq1,
\]
记录闭点的有限阶邻域；\(A_1\) 是剩余域，\(A_2\) 已能看见一次无穷小方向。与第五章的双数和切锥模型相比，这里保留整个截断环的乘法，也保留相邻阶之间的取商映射。

\begin{proposition}\label{g5:prop:formal-solutions-inverse-limit}
	设 \((R,\mathfrak n)\) 为完备且分离的局部环，
	\(F_1,\ldots,F_s\in R[Y_1,\ldots,Y_r]\)。记
	\[
		X(R/\mathfrak n^n)
		=\{y\in(R/\mathfrak n^n)^r:F_i(y)=0\text{ 对全部 }i\},
	\]
	各集合之间的映射由取模诱导。则自然映射给出双射
	\[
		X(R)\cong\varprojlim_{n\geq1}X(R/\mathfrak n^n).
	\]
\end{proposition}
\begin{proof}
	一个 \(R\) 中的解逐阶取模，自然给出相容解族。反过来，给定相容的 \(r\) 元组序列，对每个坐标使用 \(R\cong\varprojlim R/\mathfrak n^n\)，得到唯一的 \(y\in R^r\)。多项式只含有限次加法和乘法，故 \(F_i(y)\) 模每个 \(\mathfrak n^n\) 都为零。分离性给出 \(F_i(y)=0\)，于是得到解。唯一性同样逐个坐标来自分离性。
\end{proof}

相容解族要求第 \(n+1\) 层的解投影为已经选定的第 \(n\) 层解。写出各阶方程时，必须保留这些投影；单独的一层解还可能在下一层受到阻碍。例如在 \(k[[x]]\) 上，\(Y^2=x\) 模 \(x\) 有解 \(Y=0\)。若它能提升到模 \(x^2\)，提升必为 \(Y=cx\)，但 \(Y^2=0\pmod{x^2}\)，不等于 \(x\)。这个一次障碍已经排除了形式解。

\ProseParagraphBreak

对于线性方程，有限模的闭性给出更强的检测：各阶近似解不必事先相容，只要右端来自原有限模，便可直接判断它是否属于实际像。

\begin{proposition}\label{g5:prop:linear-approximation-closed-image}
	设 \((A,\mathfrak m)\) 为 Noether 局部环，\(\varphi:M\to N\) 为有限 \(A\) 模之间的同态，\(b\in N\)。若对每个 \(n\geq1\)，存在 \(u_n\in M\) 使
	\[
		b-\varphi(u_n)\in\mathfrak m^nN,
	\]
	则存在 \(u\in M\) 使 \(\varphi(u)=b\)。
\end{proposition}
\begin{proof}
	商模 \(C=N/\operatorname{im}\varphi\) 有限。条件恰好说明 \(b\) 的类属于每个 \(\mathfrak m^nC\)，故由\cref{g5:thm:krull-intersection}为零。这等价于 \(b\in\operatorname{im}\varphi\)，得到结论。
\end{proof}

这里使用的是有限模之像的闭性。若把方程换成非线性方程，便不能直接形成同一个商模；上一节的 \(Y^2=1+x\) 在 \(\mathbb Q[x]_{(x)}\) 上就有任意高阶多项式近似，却没有该环中的精确根。其形式根则在完备化中存在。

\subsection{多项式、代数幂级数与任意形式级数}
\label{g5:subsec:0705:02}

\begin{definition}\label{g5:def:algebraic-power-series}
	设 \(k\) 为域，\(x=(x_1,\ldots,x_d)\)。若 \(f\in k[[x]]\) 在有理函数域 \(k(x_1,\ldots,x_d)\) 上代数，称 \(f\) 为\term[daishumijishu]{代数幂级数}[algebraic power series]。等价地，存在非零多项式 \(P(x,Y)\in k[x,Y]\)，使
	\[
		P(x,f(x))=0.
	\]
	后一等价性来自清除有限多个系数分母。全部代数幂级数记为 \(k\langle x\rangle\)。
\end{definition}

有限多个代数元素生成有限扩域，所以它们的和、积仍代数，常数项非零的元素的形式倒数也仍代数。因此 \(k\langle x\rangle\) 是 \(k[[x]]\) 的局部子环，剩余域为 \(k\)。一元时，它就是上一节构造的环 \(S\) 的一般基域版本；下面的例子说明，多项式以外仍有许多受一个有限方程约束的级数。

\begin{example}\label{g5:exm:algebraic-square-root}
	在 \(\mathbb Q[[x]]\) 中，方程 \(Y^2=1+x\) 的常数项为 \(1\) 的根写为
	\[
		f=1+c_1x+c_2x^2+\cdots.
	\]
	比较系数得到
	\[
		2c_1=1,\qquad
		2c_n+\sum_{i=1}^{n-1}c_ic_{n-i}=0\quad(n\geq2).
	\]
	所以
	\[
		f=1+\frac{x}{2}-\frac{x^2}{8}+\frac{x^3}{16}
		-\frac{5x^4}{128}+\cdots.
	\]
	递推唯一确定全部系数，每次只需前面有限项。这是一个代数幂级数。它不是有理函数，由上一节对因子 \(x+1\) 的阶的比较即可判定，因而也不是多项式。形式解的无穷展开与定义它的一个有限方程，在这里同时存在。
\end{example}

一般形式级数还可以超出代数幂级数的范围。为给出一个可直接验证的例子，取特征零域 \(k\)，令
\[
	E(x)=\sum_{n\geq0}\frac{x^n}{n!}\in k[[x]].
\]
形式求导逐项定义，满足 \(E'=E\)。由二项式公式比较每个系数，可得 \(E(x)^j=\sum_{n\geq0}j^n x^n/n!\)，故 \((E^j)'=jE^j\)。

\begin{proposition}\label{g5:prop:formal-exponential-transcendental}
	设 \(k\) 为特征零域。形式级数 \(E(x)=\sum_{n\geq0}x^n/n!\) 在 \(k(x)\) 上超越。
\end{proposition}
\begin{proof}
	若存在代数关系，清分母后可写为
	\[
		\sum_{j\in J}p_j(x)E(x)^j=0,
	\]
	其中 \(J\) 为有限非空的非负整数集，每个 \(p_j\in k[x]\) 非零。从全部这样的关系中选取 \(\lvert J\rvert\) 最小者。若只有一项，\(k[[x]]\) 为整环且 \(E\) 为单位，已经矛盾。

	任取 \(d\in J\)，令 \(N=\deg p_d+1\)，对关系应用形式微分算子 \((\mathrm d/\mathrm dx-d)^N\)。乘积求导公式说明，第 \(j\) 项变为
	\[
		\left(\frac{\mathrm d}{\mathrm dx}+j-d\right)^N p_j(x)\;E(x)^j.
	\]
	当 \(j=d\) 时，多项式被求导 \(N\) 次，结果为零。若 \(j\ne d\)，变换后的多项式仍非零：其最高次系数为原最高次系数的 \((j-d)^N\) 倍，特征零保证该倍数非零。于是得到项数更少的非零多项式关系，与最小性矛盾。
\end{proof}

在实数或复数系数下，这个超越形式级数对每个数值 \(x\) 都收敛，通常写作 \(\exp x\)。相反，级数
\[
	G(x)=\sum_{n\geq0}n!x^n
\]
在每个非零实数或复数 \(x\) 处都不收敛，因为相邻项绝对值之比为 \((n+1)|x|\)，最终大于二，项本身不能趋于零。两者作为形式级数都完全有意义；数值收敛性由系数的增长决定，而形式运算只需每个系数由有限和给出。

\subsection{收敛、代数化及逼近的不同问题}
\label{g5:subsec:0705:03}

指定一个形式解后，可以提出三类不同的问题。收敛性问的是，在带绝对值的基域上，级数是否在某个实际邻域内定义函数。代数化问的是，给定形式对象及其指定结构，能否由有限型代数对象取完备化得到。逼近则先给定阶数 \(N\)，再寻找与形式数据模 \(\mathfrak m^N\) 相同、属于要求类别的对象。这些问题需要保留各自的输入和所要求的结构。

\begin{example}\label{g5:exm:formal-graph-marked}
	取 \(k=\mathbb Q\)、\(g(x)=E(x)-1\in xk[[x]]\)。形式代入 \(y\mapsto g(x)\) 给出
	\[
		k[[x,y]]/(y-g(x))\cong k[[x]].
	\]
	为核对核，写 \(F(x,y)=\sum_{j\geq0}a_j(x)y^j\)，并对每项使用
	\(y^j-g^j=(y-g)\sum_{i=0}^{j-1}y^{j-1-i}g^i\)。因为 \(g\in(x)\)，每个固定总次数只收到有限多项贡献，所以求和合法，得到
	\(F(x,y)-F(x,g(x))\in(y-g)\)。这证明所述核与商同构。

	这个商环本身是 \(k[x]_{(x)}\) 的完备化，因而有代数模型。但是，保留指定平面坐标 \(x,y\) 时，没有非零多项式 \(P(x,y)\in k[x,y]\) 能在整条形式图像上消失；否则
	\(P(x,E(x)-1)=0\) 会给出 \(E\) 的代数关系，违背上个命题。抽象形式局部环与它在指定平面中的图像，因而具有不同的代数化问题。

	若只要求模 \(x^N\) 相同，取 \(g\) 的次数小于 \(N\) 的截断多项式 \(g_{<N}\) 即可。多项式曲线 \(y=g_{<N}(x)\) 逐阶逼近这条形式图像，但某一次有限截断并不等于整条图像。
\end{example}

多项式近似一个给定级数容易做到；要使近似同时满足原来的多项式方程，则必须利用方程。下面在 Jacobian 可逆时给出完整的修正过程。它将上一节单根的迭代推广到同样数目的方程与未知量。

\begin{theorem}\label{g5:thm:formal-newton-system}
	设 \((R,\mathfrak n)\) 为完备且分离的局部环，
	\(F=(F_1,\ldots,F_r)\in R[Y_1,\ldots,Y_r]^r\)，\(a_0\in R^r\)。若
	\[
		F(a_0)\in(\mathfrak n)^r,\qquad
		\det J_F(a_0)\in R^\times,
	\]
	则存在唯一的 \(a\in a_0+(\mathfrak n)^r\) 满足 \(F(a)=0\)。迭代
	\[
		a_{j+1}=a_j-J_F(a_j)^{-1}F(a_j)
	\]
	满足
	\[
		F(a_j)\in(\mathfrak n^{2^j})^r,\qquad
		a-a_j\in(\mathfrak n^{2^j})^r.
	\]
	为求解模 \(\mathfrak n^N\) 的值，只需完成满足 \(2^j\geq N\) 的这些有限次修正，并可在每一步把结果截断到该商环。
\end{theorem}
\begin{proof}
	对每个分量展开有限多项式，得到
	\[
		F(a+h)=F(a)+J_F(a)h+Q(a,h),
	\]
	其中 \(Q\) 的每一项至少含两个 \(h\) 的分量。所有 \(a_j\) 都模 \(\mathfrak n\) 等于 \(a_0\)，故其 Jacobian 行列式仍为单位，矩阵逆由伴随矩阵公式给出。

	若 \(F(a_j)\in(\mathfrak n^{2^j})^r\)，则
	\(h_j=-J_F(a_j)^{-1}F(a_j)\) 的各分量也在 \(\mathfrak n^{2^j}\) 中。线性部分相消，余项落在 \((\mathfrak n^{2^{j+1}})^r\)。这就给出归纳步及 Cauchy 性，完备性和分离性给出极限 \(a\) 与 \(F(a)=0\)，也给出所述误差阶。

	为证唯一性，若 \(a,b\) 都在同一剩余类中，逐项因式分解可写
	\[
		F(b)-F(a)=(J_F(a)+C)(b-a),
	\]
	其中 \(C\) 的各条目属于 \(\mathfrak n\)。因此 \(J_F(a)+C\) 的行列式模 \(\mathfrak n\) 仍非零，是可逆矩阵，故两根相等。最后，取商保持加法、乘法及单位的逆；所以在商环中作同样的有限次修正，结果就是原迭代的截断。
\end{proof}

这里给出了达到指定形式精度的有限步骤；若系数环的有限阶商有可执行的运算，步骤就能直接实现。定理的输出首先是完备环中的精确解及其截断，并没有把每个截断多项式变成原多项式环中的精确解。Jacobian 不可逆时，前面的 \(Y^2=x\) 已出现提升障碍；一般奇异系统的逼近需要更进一步的条件与理论，\subsecref{g5:subsec:3205:02}将讨论 Artin 逼近的相应问题。
% 纲要条目排为编号小节标题，并进入目录与书签；说明段仍保留于纲要注释中。


\begin{chaptersummary}
	Rees 环把滤过的各阶和乘法放在同一个分次环中。它的 Noether 性使诱导滤过的关系由有限阶控制，给出 Artin–Rees 等式。在理想包含于 Jacobson 根时，有限模还满足 Krull 交为零，子模因而闭合。

	\ProseParagraphBreak

	完备化的正合性需要两项输入：子模滤过互相共尾，以及可数逆系统中相容提升的存在。有限呈示随后把完备化识别为张量积；伴随分次的有限生成与逐阶消去又证明完备化环 Noether。理想张量判据使这次基环扩张平坦，局部剩余域的保持则使它忠实。

	\ProseParagraphBreak

	局部完备化保留全部有限阶截断、伴随分次、长度函数、嵌入维数和 Krull 维数，也能检测有限模的自由性和理想成员关系。它容许新的形式分支及无穷展开。形式方程的解对应相容截断族；在 Jacobian 可逆时，逐次修正使误差阶加倍。一个级数是否收敛、是否满足代数关系，以及能否在指定阶数上由某类对象逼近，仍是需要分别提出的问题。
\end{chaptersummary}

\BookExercises

\begin{exercise}\label{g5:ex:rees-two-generators}
	设 \(k\) 为域，\(A=k[x,y]\)、\(I=(x,y)\)。证明 Rees 环具有呈示
	\[
		\mathcal R_I(A)\cong
		k[x,y,U,V]/(yU-xV),
	\]
	其中 \(U\) 映到 \(xt\)、\(V\) 映到 \(yt\)，并赋予 \(U,V\) 次数一、\(x,y\) 次数零。
\end{exercise}
\begin{solution}
	所给映射满射，因为 Rees 环由 \(xt,yt\) 在 \(A\) 上生成，且 \(yU-xV\) 映到零。核是分次理想，所以只需处理关于 \(U,V\) 的齐次式
	\[
		F=\sum_{j=0}^n a_j(x,y)U^jV^{n-j}
	\]
	且 \(\sum_j a_jx^jy^{n-j}=0\)。对 \(n\) 归纳。\(n=0\) 时结论直接成立。若 \(n>0\)，将等式模去 \(x\)，得到
	\(a_0(0,y)y^n=0\)，故 \(a_0=xb\)。模去理想 \((yU-xV)\)，有
	\[
		a_0V^n=xbV^n\equiv byUV^{n-1}.
	\]
	于是 \(F\) 与一个 \(UG\) 同余，\(G\) 的 \(U,V\) 次数为 \(n-1\)。将 \(U,V\) 代为 \(x,y\)，有 \(xG(x,y)=0\)。\(A\) 为整环，所以 \(G(x,y)=0\)。归纳假设给出 \(G\in(yU-xV)\)，从而 \(F\) 也在该理想中。核恰为所述主理想。
\end{solution}

\begin{exercise}\label{g5:ex:artin-rees-shift}
	取 \(A=k[x,y]\)、\(I=(x)\)、\(M=A\)、\(N=(x^r,y)\)，其中 \(r\geq1\)。计算 \(I^nM\cap N\) 与 \(I^nN\)，并求使 Artin–Rees 等式对全部 \(n\geq c\) 成立的最小非负整数 \(c\)。
\end{exercise}
\begin{solution}
	若 \(n\geq r\)，则 \(x^nA\subseteq N\)，所以交为 \(x^nA\)。若 \(n<r\)，一个 \(x^nf(x,y)\) 属于 \(N\)，当且仅当模去 \(y\) 后的 \(x^nf(x,0)\) 被 \(x^r\) 整除，等价于 \(f\in(x^{r-n},y)\)。因此
	\[
		I^nM\cap N=
		\begin{cases}
			(x^r,x^ny),&n<r,\\
			(x^n),&n\geq r,
		\end{cases}
		\qquad
		I^nN=(x^{n+r},x^ny).
	\]
	取 \(c=r\) 时，对 \(n\geq r\)，
	\(I^{n-r}(I^rM\cap N)=(x^n)\)，等于左侧。若 \(c<r\)，令 \(n=r\)，则右侧为
	\[
		I^{r-c}(I^cM\cap N)=(x^{2r-c},x^ry),
	\]
	不含 \(x^r\)：约去公共因子 \(x^r\) 后，需有 \(1\in(x^{r-c},y)\)，这是不可能的。因此最小常数为 \(r\)。诱导滤过中较低阶的两个生成方向，到第 \(r\) 阶后统一由 \(x^r\) 控制。
\end{solution}

\begin{exercise}\label{g5:ex:six-adic-completion}
	令 \(A=\mathbb Z\)。证明 \(A\) 的 \((6)\) 进完备化为
	\(\mathbb Z_2\times\mathbb Z_3\)。判断 \(A\to\widehat A\) 是否忠实平坦，并计算理想 \((5)\) 在 \((6)\) 进拓扑中的闭包。
\end{exercise}
\begin{solution}
	因为 \(2^n,3^n\) 互素，中国剩余定理给出
	\[
		\mathbb Z/6^n\mathbb Z
		\cong\mathbb Z/2^n\mathbb Z\times\mathbb Z/3^n\mathbb Z.
	\]
	同构与降阶相容，对有限乘积逐坐标取极限得到结论。完备化平坦，但不忠实：\(5\) 在每个 \(\mathbb Z/6^n\mathbb Z\) 中可逆，且逆元由唯一性组成相容族，所以它在完备化环中为单位。因此非零模 \(\mathbb Z/5\mathbb Z\) 张量后为零。

	又因 \((5)+(6^n)=\mathbb Z\)，闭包公式给出
	\[
		\overline{(5)}=\bigcap_{n\geq1}\bigl((5)+(6^n)\bigr)=\mathbb Z.
	\]
	\(\mathbb Z\) 的 \((6)\) 进拓扑本身分离，但这个真理想不闭；对应的完备化也没有保留该商模。
\end{solution}

\begin{exercise}\label{g5:ex:nilpotent-hilbert-samuel}
	设 \(k\) 为域，
	\(A=(k[x,y]/(x^2,xy))_{(x,y)}\)，极大理想为 \(\mathfrak m\)。写出 \(\widehat A\)，计算全部
	\(H(n)=\ell_A(A/\mathfrak m^{n+1})\)，并求出非零元素 \(x\) 的零化子。比较其约化环的长度函数。
\end{exercise}
\begin{solution}
	在 \(P=k[x,y]_{(x,y)}\) 上取 \((x,y)\) 进完备化，得到 \(\widehat P=k[[x,y]]\)。把\cref{g5:prop:completion-quotient}用于有限 \(P\) 模 \(P/(x^2,xy)\)，得
	\[
		\widehat A\cong k[[x,y]]/(x^2,xy).
	\]
	方程理想齐次，伴随分次环为 \(k[X,Y]/(X^2,XY)\)。也可直接在截断中读出基：\(n=0\) 时只有 \(1\)；\(n\geq1\) 时为
	\[
		1,\ x,\ y,\ y^2,\ldots,y^n.
	\]
	局部分母的常数项非零，在每个截断内已可逆，所以这些基也计算局部环的长度。因此
	\[
		H(0)=1,\qquad H(n)=n+2\quad(n\geq1).
	\]

	在完备化中，每个元素唯一写成 \(f(y)+cx\)，其中 \(f\in k[[y]]\)、\(c\in k\)。乘以 \(x\) 的结果为 \(f(0)x\)，所以 \(x\ne0\) 且其零化子为 \((x,y)\widehat A\)。由忠实平坦收缩，在 \(A\) 中也有 \(\operatorname{Ann}_A(x)=\mathfrak m\)。约化环为 \(k[y]_{(y)}\)，其长度函数为 \(n+1\)。原环多出的幂零方向只在闭点处保留一份额外长度；它没有改变最终增长的一次项，却改变了常数项。
\end{solution}

\begin{exercise}\label{g5:ex:restricted-product-completion}
	设 \(A=k[[t]]\)，\(M=\bigoplus_{j\geq0}Ae_j\)。证明其 \(t\) 进完备化可识别为
	\[
		\left\{(a_j)_{j\geq0}\in\prod_{j\geq0}A:
		\text{对每个 }n,\text{除有限多个 }j\text{ 外均有 }a_j\in t^nA\right\}.
	\]
	说明 \(A\) 完备为何仍不足以使这个自由模完备。
\end{exercise}
\begin{solution}
	给定相容族 \(u_n\in M/t^nM=\bigoplus_j(A/t^nA)e_j\)，固定坐标 \(j\) 后得到 \(A\) 的一个相容截断族，从而唯一确定 \(a_j\in A\)。每个 \(u_n\) 只有有限支撑，正好意味着除有限多个 \(j\) 外，\(a_j\) 模 \(t^n\) 为零。

	反过来，满足该条件的族在每个模 \(t^n\) 的商中只有有限支撑，给出 \(u_n\)，而这些截断自动相容。两种构造互逆并保持模运算。自然映射 \(M\to\widehat M\) 的像是只有有限多个非零坐标的族；例如 \(a_j=t^j\) 满足所述条件却不在此像中。因此 \(M\) 不完备。对有限模的完备性结论需要保留有限生成假设。
\end{solution}

\begin{exercise}\label{g5:ex:formal-node-prime-contractions}
	取
	\[
		A=\bigl(\mathbb Q[x,y]/(y^2-x^2(1+x))\bigr)_{(x,y)}.
	\]
	令 \(u\in\mathbb Q[[x]]\) 满足 \(u^2=1+x\)、\(u(0)=1\)。证明 \(\widehat A\) 的两个素理想
	\[
		\mathfrak q_\pm=(y\mp xu)\widehat A
	\]
	不同，但它们在 \(A\) 中的收缩都为零。
\end{exercise}
\begin{solution}
	商完备化给出
	\(\widehat A=\mathbb Q[[x,y]]/(y^2-x^2(1+x))\)。再模去 \(y\mp xu\)，得到 \(\mathbb Q[[x]]\)，故两个理想都为素理想。它们不同，因为在模去 \(y-xu\) 的商环中，\(y+xu\) 的像为 \(2xu\ne0\)。

	核对收缩，只需证明代入 \(y=\pm xu\) 的映射 \(A\to\mathbb Q[[x]]\) 单射。先对多项式按关于 \(y\) 的首一二次关系作除法，余式唯一为
	\(a(x)+yb(x)\)。若代入后为零且 \(b\ne0\)，就有
	\(u=\mp a/(xb)\in\mathbb Q(x)\)，与 \(1+x\) 在因子 \(x+1\) 上的阶为奇数矛盾。因此 \(b=0\)，继而 \(a=0\)。原多项式核正是方程理想。局部化的分母在原点有非零常数项，代入后为单位，所以局部化仍单射。这证明两项收缩都是零，也展示了忠实平坦谱映射可以有不同的点落在同一点上。
\end{solution}

\begin{exercise}\label{g5:ex:dvr-completion}
	设 \(A\) 为 DVR，统一化元为 \(\pi\)。证明其极大理想进完备化 \(\widehat A\) 仍为 DVR，并且 \(\pi\) 仍是统一化元，剩余域不变。
\end{exercise}
\begin{solution}
	\(\widehat A\) 为 Noether 局部环，极大理想为 \(\pi\widehat A\)，剩余域由\cref{g5:thm:completion-faithful}保持。该定理还保证 \(A\hookrightarrow\widehat A\)，所以 \(\pi\) 及其全部幂均非零。

	对 \(0\ne b\in\widehat A\)，分离性给出唯一的非负整数 \(n\)，使
	\(b\in\pi^n\widehat A\setminus\pi^{n+1}\widehat A\)。写 \(b=\pi^n v\)，则 \(v\notin\pi\widehat A\)，故 \(v\) 为单位。两个非零元素的乘积因而为 \(\pi^{n+m}\) 乘以单位，仍非零，说明 \(\widehat A\) 为整环。

	任一非零理想中，元素的这些阶有最小值 \(r\)。一个达到最小阶的元素生成 \(\pi^r\widehat A\)，而全部其他元素也属于该主理想，所以原理想等于它。由此得到全部非零理想为统一化元的幂，正是 DVR 的刻画；极大理想非零且真，故该环不是域。
\end{solution}

\begin{exercise}\label{g5:ex:localization-of-complete-ring}
	设 \(R=k[[x,y]]\)、\(A=R_{(x)}\)。证明
	\[
		\widehat A\cong k((y))[[x]],
	\]
	其中左侧取 \(x\) 进完备化。再证明 \(A\) 不完备。可考虑级数
	\(h=\sum_{n\geq0}y^{-n!}x^n\)。
\end{exercise}
\begin{solution}
	模去 \(x^N\) 后，每个级数由有限多个 \(x\) 次幂及 \(k[[y]]\) 系数表示。\(R\setminus(x)\) 中的分母，其 \(x^0\) 系数是非零的 \(k[[y]]\) 元素；在 \(k((y))[x]/(x^N)\) 中，先除去这个常数项，再用有限几何级数求逆。因此
	\[
		A/x^NA\cong k((y))[x]/(x^N).
	\]
	满射也可直接看出：右侧有限多个 Laurent 系数有共同的 \(y\) 幂分母，而 \(y\notin(x)\)。这些同构与降阶相容，取极限得到完备化公式。

	若 \(h=f/g\in A\)，将 \(f,g\in R\) 按 \(x\) 展开，记 \(g_0(y)\ne0\)，\(d=\operatorname{ord}_y g_0\geq0\)。一般地，若
	\(f/g=\sum b_n(y)x^n\)，系数递推
	\[
		g_0b_n=f_n-\sum_{j=1}^n g_jb_{n-j}
	\]
	给出 \(\operatorname{ord}_y b_n\geq-d(n+1)\)。确实，\(f_n,g_j\in k[[y]]\) 的阶非负；从 \(n=0\) 起归纳，每步除以 \(g_0\) 至多再损失 \(d\) 阶，便得到这个下界。

	但 \(h\) 的第 \(n\) 个系数阶为 \(-n!\)，对充分大的 \(n\) 小于任何固定的 \(-d(n+1)\)，矛盾。因此 \(h\notin A\)。原来的 \(R\) 完备，其局部化 \(A\) 却不完备。
\end{solution}

\begin{exercise}\label{g5:ex:idempotent-lifting-complete}
	设 \(R\) 为交换含幺环，对理想 \(I\) 的进拓扑完备且分离。证明每个 \(R/I\) 中的幂等元都唯一提升为 \(R\) 中的幂等元。
\end{exercise}
\begin{solution}
	首先，若 \(u\in I\)，几何级数 \(\sum_{n\geq0}u^n\) 给出 \(1-u\) 的逆。因此 \(I\subseteq\operatorname{Jac}(R)\)。选所给幂等元的提升 \(e_0\)，则 \(e_0^2-e_0\in I\)，并且
	\[
		(2e_0-1)^2=1+4(e_0^2-e_0)
	\]
	为单位，故 \(2e_0-1\) 也为单位。这一论证对任意与 \(e_0\) 模 \(I\) 同余的元素成立，无需假设 \(2\) 可逆。

	对 \(F(T)=T^2-T\) 采用单根证明中的迭代
	\[
		e_{j+1}=e_j-(e_j^2-e_j)(2e_j-1)^{-1}.
	\]
	同一多项式展开给出误差从 \(I^{2^j}\) 进入 \(I^{2^{j+1}}\)，所以迭代收敛到幂等元 \(e\)，并满足 \(e\equiv e_0\pmod I\)。

	若另一幂等元 \(f\equiv e\pmod I\)，则
	\[
		0=(f^2-f)-(e^2-e)=(f-e)(f+e-1).
	\]
	第二因子模 \(I\) 等于 \(2e_0-1\)，故为单位，推出 \(f=e\)。
\end{solution}

\begin{exercise}\label{g5:ex:hensel-characteristic-two}
	在 \(\mathbb F_2[[x]]\) 中，证明 \(Y^2=1+x\) 无解，而 \(Y^2+Y=x\) 有唯一常数项为零的解，并将它写为显式形式级数。
\end{exercise}
\begin{solution}
	特征二时，一个级数的平方只含偶次幂，所以不可能等于含 \(x\) 项的 \(1+x\)。第二个方程的导数为 \(1\)，Hensel 单根提升保证常数项零的解唯一。令
	\[
		y=\sum_{j\geq0}x^{2^j}=x+x^2+x^4+x^8+\cdots.
	\]
	逐阶平方得到 \(y^2=\sum_{j\geq0}x^{2^{j+1}}\)，所以 \(y^2+y=x\)，给出所求解。它由一个二次多项式约束，是代数幂级数。两个方程的单根条件不同，决定了从剩余域根出发是否能进行提升。
\end{solution}

\begin{exercise}\label{g5:ex:newton-square-root-precision}
	对 \(\mathbb Q[[x]]\) 中的 \(Y^2=1+x\)，以 \(a_0=1\) 开始作 Newton 修正。说明三次修正足以得到模 \(x^8\) 的根，并写出这一级精度的多项式。
\end{exercise}
\begin{solution}
	误差 \(F(a_0)=-x\) 从一阶开始；单根定理使第 \(j\) 次修正后的误差属于 \((x^{2^j})\)，且当前值与精确根之差也属于这个理想。因此三次修正足够。前两步为
	\[
		a_1=1+\frac{x}{2},\qquad
		a_2=\frac{x^2+8x+8}{4x+8}.
	\]
	每个分母在形式环中为单位。第三步按
	\(a_3=\tfrac12(a_2+(1+x)/a_2)\) 计算并截断，或使用正文的系数递推，得到
	\[
		a_3\equiv
		1+\frac{x}{2}-\frac{x^2}{8}+\frac{x^3}{16}
		-\frac{5x^4}{128}+\frac{7x^5}{256}
		-\frac{21x^6}{1024}+\frac{33x^7}{2048}
		\pmod{x^8}.
	\]
	将此多项式平方，零次项为一、一次项为 \(x\)，二至七次项均相消。这是指定精度的解；作为多项式，它的平方仍可能具有八次及以上的非零误差。
\end{solution}

\begin{exercise}\label{g5:ex:formal-coordinate-inverse}
	设 \(k\) 为域，\(f_1,\ldots,f_d\in(x_1,\ldots,x_d)k[[x_1,\ldots,x_d]]\)。若其一次项组成可逆矩阵 \(L\in\operatorname{GL}_d(k)\)，证明代入 \(x_i\mapsto f_i\) 给出形式幂级数环的连续自同构。
\end{exercise}
\begin{solution}
	各 \(f_i\) 常数项为零，所以在任意固定总次数内，代入只涉及原级数的有限多项；代入有定义，保持环运算且连续。记向量 \(f=Lx+\) 高次项。我们逐次构造 \(g\in(x)^d\)，使 \(f(g)=x\)。一次项必须为 \(L^{-1}x\)。若已选次数小于 \(N\) 的 \(g\)，且误差 \(f(g)-x\) 属于 \((x)^N\)，记其次数 \(N\) 的齐次部分为 \(r_N\)。添加齐次修正
	\[
		h_N=-L^{-1}r_N
	\]
	后，\(f(g+h_N)-f(g)\) 的次数 \(N\) 部分恰为 \(Lh_N\)，其余新增项次数至少为 \(N+1\)。因此消去该阶误差，递归给出唯一的 \(g\)，满足 \(f(g)=x\)。

	最后验证双侧互逆。任何一次项矩阵可逆的代入都单射：非零级数的最低次齐次部分在可逆线性换元后仍非零。记两张代入映射为 \(\varphi_f,\varphi_g\)，已得
	\(\varphi_g\varphi_f=1\)。由于 \(g\) 的一次项矩阵为 \(L^{-1}\)，\(\varphi_g\) 也单射。对
	\(\varphi_g(\varphi_f\varphi_g-1)=0\) 使用这一单射性，得
	\(\varphi_f\varphi_g=1\)。两张代入均连续，因此是所求拓扑环自同构。
\end{solution}
